% Options for packages loaded elsewhere
\PassOptionsToPackage{unicode}{hyperref}
\PassOptionsToPackage{hyphens}{url}
% custom: no * to svgnames, causes some colors not to be recognized
\PassOptionsToPackage{dvipsnames,svgnames,x11names*}{xcolor}
%
\documentclass[american,a4paper,]{article}
\usepackage{lmodern}
\usepackage{texgyretermes}
\usepackage{amssymb,amsmath}
\usepackage{ifxetex,ifluatex}
\usepackage{fixltx2e} % provides \textsubscript


% use upquote if available, for straight quotes in verbatim environments
\IfFileExists{upquote.sty}{\usepackage{upquote}}{}
% use microtype if available
\IfFileExists{microtype.sty}{%
\usepackage[]{microtype}
\UseMicrotypeSet[protrusion]{basicmath} % disable protrusion for tt fonts
}{}
\IfFileExists{parskip.sty}{%
\usepackage{parskip}
}{% else
\setlength{\parindent}{0pt}
\setlength{\parskip}{6pt plus 2pt minus 1pt}
}
% custom: footnotes (more below)
\usepackage[bottom,hang,multiple,stable]{footmisc}
\setlength{\footnotesep}{4mm}
\usepackage{xcolor}
\usepackage{hyperref}
\hypersetup{
            pdftitle={Intelligenza artificiale e funzione giurisdizionale},
 % custom
              pdfauthor={Giuliano Taglianetti, Raoul Amitrano},
              pdflang={en-US},
            colorlinks=true,
            linkcolor=DarkBlue,
            citecolor=DarkBlue,
            urlcolor=DarkBlue,
            breaklinks=true}
\urlstyle{same}  % don't use monospace font for urls
\usepackage[a4paper,headheight=20pt,bottom=100pt,left=1.25in,right=1.25in]{geometry}
% custom: to improve separations
\widowpenalty=1000
\clubpenalty=1000
% custom: invoke sectsty, we are going to use it 95% of the time
\usepackage{sectsty}
% custom: Title Font for headings and titles (\titlefont is set above)
\usepackage{graphicx}
%\usepackage[normalem]{ulem}
\setlength{\emergencystretch}{3em}  % prevent overfull lines
\providecommand{\tightlist}{%
  \setlength{\itemsep}{0pt}\setlength{\parskip}{0pt}}
\setcounter{secnumdepth}{5}
% Redefines (sub)paragraphs to behave more like sections
\ifx\paragraph\undefined\else
\let\oldparagraph\paragraph
\renewcommand{\paragraph}[1]{\oldparagraph{#1}\mbox{}}
\fi
\ifx\subparagraph\undefined\else
\let\oldsubparagraph\subparagraph
\renewcommand{\subparagraph}[1]{\oldsubparagraph{#1}\mbox{}}
\fi

% custom: suppress numbering in figures and tables, caption centered and small
\usepackage[font=small,justification=centering]{caption}
\captionsetup[figure]{labelformat=empty}
\captionsetup[table]{labelformat=empty}

% Workaround/bugfix from jannick0.
% See https://github.com/jgm/pandoc/issues/4302#issuecomment-360669013)
% or https://github.com/Wandmalfarbe/pandoc-latex-template/issues/2
%
% Redefine the verbatim environment 'Highlighting' to break long lines (with
% the help of fvextra). Redefinition is necessary because it is unlikely that
% pandoc includes fvextra in the default template.
\usepackage{fvextra}
\DefineVerbatimEnvironment{Highlighting}{Verbatim}{breaklines,fontsize=\footnotesize,commandchars=\\\{\}}

\renewenvironment{quote}%
               {\list{}{\rightmargin\leftmargin}%
                \item[]\samepage}% <- The effect of \samepage is local!!!
               {\endlist}

% custom: icons
% Both fontawesome5 and academicons are injected into the compilation by pipeline.js
\usepackage{fontawesome5}
\usepackage{academicons}

% custom: dates
\newcommand{\DTMdate}[1]{#1}%

% custom: page number
  \setcounter{page}{1}

% custom: footer and header (except first page)
\usepackage[breakwords]{truncate}
\usepackage{fancyhdr}
\pagestyle{fancy}
\fancyhf{} % clear all default
% left header
\lhead{\footnotesize\sffamily \truncate{9cm}{Intelligenza artificiale e funzione giurisdizionale}}
% right reader
\rhead{\footnotesize\sffamily i-lex. Vol. 19 n.~1 (2026)}
% left footer (DOI)
\lfoot{\footnotesize\sffamily \url{https://doi.org/10.60923/issn.1825-1927/22648}}
% right footer (ID (if issuepagination false) + page (with optional prefix))
\rfoot{\footnotesize\sffamily  \thepage}
% custom: TEST for more behaved footnotes
\usepackage{etoolbox}
% no superscript footnote marks, by egreg from https://tex.stackexchange.com/a/96400
\makeatletter
\newcommand{\nosupmark@makefnmark}{\@thefnmark.}
\pretocmd{\@makefntext}{\let\@makefnmark\nosupmark@makefnmark}{}{}

% custom: footnotes in first page similar to normal ones
\patchcmd\maketitle{\renewcommand\thefootnote{\@alph\c@footnote}}{\AdaptNote\thanks\multthanks}{}{}
\patchcmd\maketitle{%
  \def\@makefnmark{\rlap{\@textsuperscript{\normalfont\@thefnmark}}}}{}{}{}
\makeatother
% custom: hideFromPandoc command to allow latex block with \Begin and \End
\newcommand{\hideFromPandoc}[1]{#1}
\hideFromPandoc{
  \let\Begin\begin
  \let\End\end
}
%\usepackage{soul}
% custom: define \linethickness
\renewcommand{\linethickness}{0.3pt}
% custom: draft mark with date

% set default figure placement to htbp
\makeatletter
\def\fps@figure{htbp}
\makeatother

\ifnum 0\ifxetex 1\fi\ifluatex 1\fi=0 % if pdftex
  \usepackage[shorthands=off,main=american]{babel}
\fi
% update 2020

\title{Intelligenza artificiale e funzione giurisdizionale\\ \large Profili applicativi nel processo amministrativo}

\renewcommand{\thefootnote}{\fnsymbol{footnote}}

 % custom: affiliation only if "affiliationonly" is true, else boxed version below abstract
\author{
{\small Giuliano Taglianetti}
{\small\textsuperscript{1}}
 \and {\small Raoul Amitrano}
{\small\textsuperscript{1}}
}


\renewcommand{\thefootnote}{\arabic{footnote}}

% custom: dates


\date{}

 % end custom dates

%\setlength{\thanksmarkwidth}{1.8em}
%\thanksfootextra{}{\hspace{5pt}}



\begin{document}
\maketitle
\fancypagestyle{plain}{%
  \fancyhf{}%
  % right header: journal title (extended) and eISSN
  \rhead{\footnotesize\sffamily i-lex -- Rivista di Scienze Giuridiche, Scienze Cognitive ed Intelligenza Artificiale. Vol. 19 n.~1 (2026) \\ ISSN~1825-1927}
  % left header: section, peer-reviewed plus DOI or journal url
    \lhead{\footnotesize\sffamily Essays -- \emph{peer-reviewed}
   \\ \url{https://doi.org/10.60923/issn.1825-1927/22648}}
    % right footer (ID + page (with optional prefix))
  \rfoot{\footnotesize\sffamily  \thepage}
    % left footer: copyright block
  \lfoot{
  \scriptsize\sffamily
  Copyright \copyright~2026 Giuliano Taglianetti, Raoul Amitrano \\ This work is licensed under the Creative Commons BY License. \\ \url{https://creativecommons.org/licenses/by/4.0/}
  }
  }

\begin{figure}[!b]
\begin{minipage}{0.9\textwidth}
\scriptsize\sffamily
\faEnvelope
~example@email.org (Giuliano Taglianetti);
~example@email.org (Raoul Amitrano);
\\
\aiOrcid~\url{https://orcid.org/0000-0002-1825-0097} (Giuliano Taglianetti);
\end{minipage}
\end{figure}

\newenvironment*{institute}{}{}
\begin{institute}
\footnotesize
\textsuperscript{1} University of Napoli Federico II, Italy\\
\end{institute}



% custom: no author blocks, just affiliation
 % if(title)
 % custom: to opt-out to toc in pdf
 % custom: - end if(nopdftoc)
% custom: space from title or separate title page
\vspace{1cm}
\textbf{Abstract}: Il contributo analizza l'impatto dell\textquotesingle intelligenza artificiale sulla funzione giurisdizionale, con particolare riferimento al processo amministrativo e all'esercizio della discrezionalità giudiziale. Muovendo dal progressivo ingresso di strumenti algoritmici nelle attività decisionali pubbliche, l'indagine esamina il ruolo dell'IA quale supporto all'attività del giudice e ne valuta i limiti strutturali rispetto alle funzioni valutative e di bilanciamento proprie della giurisdizione amministrativa.

\textbf{Keywords}: Intelligenza artificiale; processo amministrativo; discrezionalità giudiziale; algoritmo; decisione giudiziale; giudice amministrativo; bilanciamento degli interessi; processo cautelare; contratti pubblici; appalti; effettività della tutela giurisdizionale; automatizzazione decisionale.

\section{Introduzione: premessa, oggetto e scopo dell'indagine}\label{introduzione-premessa-oggetto-e-scopo-dellindagine}

La progressiva integrazione di sistemi algoritmici e di intelligenza artificiale nei processi decisionali pubblici rappresenta uno dei fenomeni più significativi della trasformazione contemporanea del diritto amministrativo, in particolare nei settori connotati da un elevato grado di discrezionalità. Ciò che, fino a tempi recenti, si configurava come una prospettiva emergente ha ormai assunto i tratti di un elemento strutturale dell'ordinamento, imponendo a dottrina e giurisprudenza un ripensamento critico delle categorie tradizionali attraverso cui è stato storicamente definito il rapporto tra amministrazione, giudice e cittadino. In questo contesto, il processo amministrativo si configura come un osservatorio privilegiato delle tensioni generate dall'interazione tra intelligenza artificiale e discrezionalità. Esso non si esaurisce necessariamente nella mera attività ermeneutica e nella verifica della legittimità dell'azione amministrativa, ma può divenire la sede in cui si realizza un delicato bilanciamento tra la tutela delle posizioni giuridiche soggettive e la salvaguardia dell'interesse pubblico sotteso al provvedimento oggetto di giudizio. In tale quadro, il giudice amministrativo è frequentemente chiamato a esercitare poteri che travalicano la sola interpretazione della norma, implicando valutazioni complesse e un'attenta ponderazione tra interessi pubblici e privati eterogenei. È proprio in queste ipotesi che si manifesta la c.d. discrezionalità giudiziale, intesa quale attività valutativa che integra elementi normativi, fattuali e assiologici e che non può essere ridotta a una mera applicazione automatica di regole predeterminate\footnote{* Giuliano Taglianetti è Professore associato di Diritto amministrativo presso l'Università degli Studi di Napoli ``Federico II''; Raoul Amitrano è laureato in giurisprudenza presso il medesimo Ateneo. Il contributo è frutto di elaborazione comune. In particolare, i paragrafi nn.~1 e 5 sono attribuibili a Giuliano Taglianetti, mentre i paragrafi nn.~2, 3 e 4 sono attribuibili a Raoul Amitrano; le conclusioni sono state redatte congiuntamente.

  Per un'approfondita analisi della discrezionalità nel campo giudiziario, cfr. F. Bricola, \emph{La discrezionalità nel diritto penale}, Milano, 1965; A. Raselli, \emph{Studi sul potere discrezionale del giudice civile}, Milano, 1975; A. Barak, \emph{JudicialDiscrection}, New Haven, 1989, trad. it. a cura di I. Mattei, \emph{La discrezionalità del giudice}, Milano, 1995; R. Guastini, \emph{Principi di diritto e discrezionalità giudiziale}, in \emph{Dir. pubbl}., 1998, p.~641 ss. Per una disamina della discrezionalità del giudice amministrativo cfr., P. Gotti,\emph{Considerazioni su sorte del contratto di appalto pubblico e discrezionalità del giudice amministrativo}, in \emph{Ann. Fac. Giur. Univ. Camerino}, n.~1/2012; F. Saitta, \emph{Interprete senza spartito? Saggio critico sulla discrezionalità del giudice amministrativo}, Napoli, 2023; G. Taglianetti, \emph{Il processo amministrativo tra giustizia e certezza. Le regole dell'equilibrio tra l'effettività della tutela giurisdizionale e la stabilità delle decisioni pubbliche,} Milano, 2023, spec. p.~153 ss., a cui si rinvia anche per ulteriori e ampi richiami bibliografici di dottrina.}. Un ambito emblematico di tale fenomeno è rappresentato dal procedimento cautelare, nel quale il giudice è chiamato a svolgere una valutazione comparativa tra gli interessi individuali delle parti e quelli pubblici coinvolti, soppesandone la qualità e la rilevanza alla luce di un prudente e motivato apprezzamento. Particolarmente significative, al riguardo, sono le decisioni adottate nel periodo dell'emergenza Covid-19: nei giudizi aventi ad oggetto i provvedimenti di chiusura delle scuole, il giudice amministrativo è stato chiamato a svolgere complesse valutazioni di bilanciamento tra diritti fondamentali, quali, da un lato, il diritto all'istruzione dei minori e il diritto al lavoro dei genitori e, dall'altro, il diritto alla salute, anche nella sua dimensione collettiva, in quanto connesso all'esigenza di prevenire la diffusione del virus\footnote{In tal senso, si veda, tra le altre, Cons. Stato, sez. III, 10 novembre 2020, n.~6453, in \emph{Giustiziainsieme}, 2020, con nota di C. Napolitano, \emph{Regioni, scuola e COVID-19: il Giudice Amministrativo tra diritto allo studio e tutela della salute}, nella quale il Collegio fa espresso riferimento a una «ponderazione comparata degli interessi in contrasto». In argomento, v. R. Dagostino, \emph{Emergenza pandemica e tutela cautelare (monocratica)}, in \emph{Giustiziainsieme.it}, 2020; F. Saitta,~\emph{Sulla decisione di prevedere una tutela cautelare monocratica}~ex officio~\emph{nell'emergenza epidemiologica da Covid-19: chi? Come? ma soprattutto, perché?},~in \emph{federalismi.it}, 2020;~Id.,~\emph{Il processo cautelare alle prese con la pandemia}, in \emph{\href{http://www.lexitalia.it}{www.lexitalia.it},} 2020; C. Cacciavillani,~\emph{Controcanto sulla disciplina emergenziale del processo amministrativo (con riferimento all'art. 4 del d.l. 30 aprile 2020, n.~28)}, in~\emph{GiustAmm}, n.~5/2020, p.~1 ss.;~F. Francario,~\emph{Diritto dell'emergenza e giustizia nell'amministrazione. No a false semplificazioni e a false riforme}, in~\emph{Federalismi}, 2020, p.~3 ss.; N. Paolantonio,~\emph{Il processo amministrativo dell'emergenza: sempre più speciale}, in~\emph{GiustAmm}, 2020, 1 ss.}. Ne deriva che l'esercizio della discrezionalità cautelare non può esaurirsi in una valutazione incentrata esclusivamente sulla posizione soggettiva del ricorrente, ma richiede un bilanciamento che consideri anche l'incidenza della misura richiesta sull'efficace perseguimento delle finalità pubbliche. Ciò in quanto manca una scala gerarchica, e persino una predeterminazione legislativa, degli interessi coinvolti, i quali sono invece ricavabili dai valori dell'ordinamento e dai principi che lo informano. A ciò si aggiunge la scelta tra diverse tipologie di misure cautelari - sospensive, propulsive o anticipatorie, talora anche con effetti irreversibili - scelta intrinsecamente connessa all'individuazione dell'interesse ritenuto, nel caso concreto, meritevole di prevalenza. Dinamiche analoghe, sebbene non integralmente sovrapponibili - come si avrà modo di evidenziare nel prosieguo - emergono altresì nel contenzioso in materia di contratti pubblici, nell'ambito del quale il giudice amministrativo è chiamato a contemperare l'interesse del ricorrente a una tutela piena ed effettiva con l'interesse pubblico alla celere realizzazione dell'opera ovvero alla continuità del servizio. In particolare, nelle ipotesi previste dall'art. 122 del decreto legislativo 2 luglio 2010, n.~104 (Codice del processo amministrativo), tale discrezionalità si caratterizza per l'assenza di una predeterminazione legislativa di una gerarchia tra gli interessi coinvolti, essendo rimessa al giudice l'individuazione, caso per caso, del punto di equilibrio più adeguato. Risulta evidente l'ampiezza del potere di scelta riconosciuto al giudice amministrativo, inteso quale autentico ``spazio di manovra'' normativamente attribuito. Si tratta di un ambito decisionale che, lungi dall'essere rigidamente predeterminato, implica l'esercizio di una valutazione orientata all'individuazione, nel caso concreto, del punto di equilibrio più adeguato tra gli interessi in gioco, secondo criteri di ragionevolezza e proporzionalità. Tali valutazioni costituiscono espressione della discrezionalità giudiziale e devono essere condotte alla luce del principio di effettività della tutela giurisdizionale, cardine del processo amministrativo, espressamente enunciato all'art. 1 c.p.a., la cui collocazione sistematica ne evidenzia la centralità nell'architettura complessiva del processo. Esse devono, tuttavia, confrontarsi anche con l'interesse pubblico concreto perseguito dal provvedimento impugnato e, quindi, con il principio del risultato, desumibile in via implicita da numerose disposizioni del codice del processo amministrativo, quale specifica declinazione del principio costituzionale di buon andamento dell'azione amministrativa\footnote{Cfr. G. Taglianetti, \emph{Processo amministrativo e principio del risultato}, in \emph{Dir. proc. amm.}, 2025, p.~510 ss.}. La discrezionalità giudiziale si distingue, dunque, in modo sottile ma significativo, dalla discrezionalità amministrativa, potendo giungere a dare attuazione a norme generali prive di una fattispecie puntuale, quali quelle di rango costituzionale, nonché a operare il contemperamento tra principi e valori, oltre che tra istanze individuali e interessi collettivi. È proprio all'interno di questo spazio valutativo che l'introduzione di strumenti di intelligenza artificiale solleva le questioni teoriche più delicate. Da un lato, tali strumenti promettono rilevanti benefici in termini di efficienza, coerenza decisionale e gestione della complessità, particolarmente apprezzabili in un ambito - quale quello amministrativo - in cui il giudice mostra talora una marcata propensione all'elaborazione creativa e all'interpretazione per principi\footnote{Il tema della creatività giurisprudenziale - e, più in generale, il rapporto tra giudice e legislatore - costituisce un ambito di indagine particolarmente ampio e stratificato, da tempo al centro di riflessioni dottrinali anche profondamente divergenti. Senza alcuna pretesa di esaustività, possono richiamarsi, tra i numerosi contributi in materia: H. Kelsen,~\emph{Sulla teoria dell'interpretazione}~(1934), in Id.,~\emph{Il primato del Parlamento}, a cura di C. Geraci, Milano, 1982, p.~159 ss.; Id., \emph{Lineamenti di dottrina pura del diritto} (1934) trad. it. a cura di R. Treves, Torino, 1970; Id., \emph{La dottrina pura del diritto} (1960), trad. it. di M. G. Losano, Torino, 1966; V. Crisafulli, \emph{I principi costituzionali dell'interpretazione ed applicazione delle leggi}, in \emph{Scritti giuridici in onore di Santi Romano}, I, Padova, 1940, pp.~665 ss.; E. Betti, \emph{Interpretazione della legge e degli atti giuridici (Teoria generale e dogmatica)}, Milano, 1949, pp.~163 ss.; M. Nigro, \emph{Il Consiglio di Stato giudice ed amministratore. Aspetti di effettività dell'organo}, in \emph{Riv. trim. dir. proc. civ}., 1974, p.~1371 ss., poi in \emph{Scritti giuridici in onore} di S. Pugliatti, II, Milano, 1978, p.~993 ss., infine in M. Nigro, \emph{Scritti giuridici}, Milano, 1996, II, p.~1051 ss.; M. Cappelletti, \emph{Giudici legislatori?}, Milano, 1984; N. Lipari, \emph{Il ruolo del giudice nella crisi delle fonti del diritto}, in \emph{Riv. trim. dir. proc. civ.}, 2009, p.~479 ss; R. Guastini,~\emph{Interpretare e argomentare}, Milano, 2011; G. Pino,~\emph{Interpretazione cognitiva, interpretazione decisoria},~\emph{interpretazione creativa}, in~\emph{Riv. fil. dir.,} 2013, p.~77 ss.; Id. \emph{Poteri interpretativi e principio di legalità}, in \emph{Enciclopedia del diritto. I tematici}, vol.~V, \emph{Potere e costituzione}, Milano, 2023, pp.~960 ss.; S. Spuntarelli, \emph{La giurisprudenza \textquotesingle creativa\textquotesingle{} del giudice amministrativo}, in P. Di Lucia, L. Passerini Glazel, \emph{Il dover essere del diritto: un dibattito teorico sul diritto illegittimo a partire da Kelsen}, Torino 2020, p.~241 ss.; M. Barberis,~\emph{Beccaria, Bentham e il creazionismo giuridico}, in~\emph{Riv. internaz. fil. dir.}, 2014, p.~539 ss.; M. Luciani, \emph{Il ruolo della giurisprudenza: applicazione delle clausole generali e certezza dei rapporti giuridici}, in G. Grasso, M. Maugeri (a cura di), \emph{Le fonti del diritto, il ruolo della giurisprudenza e il principio di legalità}, Roma, 2023, p.~171 ss.; Id., \emph{Ogni cosa al suo posto}, Milano, 2023, spec. p.~147 ss.; A. Lucarelli, \emph{Diritto e conflitti nell'incertezza della produzione giuridica}, in \emph{Rivistaaic}, 2024, p.~77 ss. Sullo stesso tema, si segnalano inoltre i contributi pubblicati in \emph{Dir. pubb.}, 2016/2, dedicato a \emph{Giudici e legislatori}, con interventi di C. Pinelli, G.U. Rescigno, A. Travi, M. Bombardelli, V. Angiolini, P. Ciarlo, A. Pioggia, S. Civitarese Matteucci, G. Azzariti, F. Bilancia, P. Carnevale, D. Sorace), e in \emph{Questione giustizia}, 2016/4, dedicato a \emph{Il giudice e la legge}, con scritti di L. Ferrajoli, N. Lipari, F. Macario, G. Preterossi, P.L. Zanchetta, R.G. Conti, A. Lamorgese, E. Scoditti, A. Giusti, A. Terzi, G. Leo, D. Simeoli, A. Guardiano.}. Dall'altro, essi sollevano interrogativi cruciali, non esclusivamente di natura giuridica\footnote{Il dibattito sull'intelligenza artificiale solleva questioni complesse, tra cui il significato stesso di intelligenza, ciò che conosciamo e ignoriamo del nostro modo di pensare, i limiti qualitativi e quantitativi del calcolo, nonché il problema dell'allineamento della tecnologia ai nostri bisogni e ai nostri valori. Con riguardo a quest'ultimo aspetto, la riflessione può essere approfondita in una prospettiva filosofico-antropologica, che invita a interrogarsi sull'effettiva auspicabilità di sostituire il giudice umano con sistemi di intelligenza artificiale, anche nell'ipotesi in cui questi ultimi fossero in grado di garantire risultati migliori in termini di giustizia sostanziale, equità delle decisioni e adeguatezza rispetto alla complessità degli interessi coinvolti. Su un piano ancora più generale, l'analisi conduce a domandarsi se possano concepirsi forme di coscienza diverse da quella umana. Tuttavia, il profilo più problematico non sembra risiedere tanto nella possibilità di una sostituzione integrale dell'essere umano da parte della macchina, quanto piuttosto nel rischio - più sottile ma pervasivo - che l'uomo finisca per identificarsi nel modello algoritmico, arrivando a concepirsi come un insieme di processi deterministici e di reazioni integralmente codificabili.}, circa la loro compatibilità con una funzione giurisdizionale fondata sulla sensibilità del giudice e sull'esercizio di un potere valutativo non integralmente riducibile a un calcolo algoritmico. In questa prospettiva, la discrezionalità giudiziale rappresenta un banco di prova privilegiato per verificare i limiti strutturali dell'automazione decisionale nel processo amministrativo. La presente indagine assume pertanto come fulcro il rapporto tra intelligenza artificiale, processo amministrativo e discrezionalità giudiziale, con l'obiettivo di accertare se, e in quale misura, l'impiego di strumenti algoritmici possa incidere sull'esercizio della funzione giurisdizionale, distinguendo tra un utilizzo meramente ausiliario dell'intelligenza artificiale e ipotesi di possibile sostituzione del decisore umano. Il punto fermo della presente analisi è che l'intelligenza artificiale e gli strumenti automatizzati hanno un'indubbia utilità come strumenti ausiliari del processo e meritano certamente di essere implementati, ad esempio per l'organizzazione delle udienze, la gestione dei ruoli o la ricerca intelligente di precedenti giurisprudenziali\footnote{Su tali aspetti, v. H. Simonetti, \emph{Brevi note sul possibile uso dell'intelligenza artificiale nel processo amministrativo}, in \emph{P.A. Persona e Amministrazione}, n.~1/2024, p.~917 ss.}. Partendo da questo presupposto, la ricerca si sviluppa su due livelli complementari. Innanzitutto, si analizza il ruolo dell'intelligenza artificiale nel processo in generale, con particolare riguardo all'attività interpretativa del giudice, anche in relazione al precedente giurisprudenziale e alla funzione di orientamento delle decisioni. In questo contesto, è fondamentale interrogarsi sulla capacità dei sistemi di IA di autoistruirsi e autoimplementarsi, potenzialmente dando origine a fenomeni analoghi ai \emph{revirement} giurisprudenziali. Allo stesso tempo, occorre considerare il rischio opposto: che l'utilizzo dell'IA, basandosi prevalentemente sui precedenti, possa appiattire il giudizio, riducendo la capacità del giudice di adeguare le decisioni alle nuove situazioni concrete. Dunque, si verifica il possibile impatto di tali strumenti sulla discrezionalità giudiziale del giudice amministrativo, distinta dall'attività interpretativa poiché implica valutazioni assiologiche complesse e scelte di priorità tra interessi pubblici e privati, come avviene nelle decisioni cautelari o in quelle relative alla sorte dei contratti in materia di appalti. L'indagine mostra come l'equilibrio tra effettività della tutela giurisdizionale e tutela dell'interesse pubblico concreto perseguito dal provvedimento impugnato costituisca il nucleo della giustizia amministrativa, rendendo la discrezionalità giudiziale uno strumento imprescindibile e il giudice il fulcro insostituibile del processo.

\section{Miti ed equivoci nella materia algoritmica}\label{miti-ed-equivoci-nella-materia-algoritmica}

Prima ancora di valutarne l'impatto sulla funzione giurisdizionale, appare necessario chiarire cosa debba intendersi per ``intelligenza artificiale''\footnote{Si rinvia, sul tema, all'opera di G. Sartor, \emph{L'intelligenza artificiale e il diritto}, Torino, 2022, spec. pp.~1-33, nonché alla dottrina ivi richiamata.}. La discussione pubblica sull'IA è spesso viziata da un equivoco concettuale: molti sistemi oggi diffusi, inclusi i modelli linguistici di grandi dimensioni, non sono propriamente ``intelligenze artificiali''\footnote{Per una visione d'insieme sul tema, si veda l'autorevole manuale di S. Russell e P. Norvig, \emph{Intelligenza artificiale. Un approccio moderno}, IV ed., Milano, Pearson, 2021; per una efficace sintesi dei principali profili tecnici, v. Pascuzzi, \emph{La cittadinanza digitale. Competenze, diritti e regole per vivere in rete}, Bologna, Il Mulino, 2021.}, ma strumenti di linguistica computazionale basati su correlazioni statistiche. Essi non comprendono la realtà sensibile, non la percepiscono, e non la interpretano: predicono. La loro utilità risiede nella capacità di individuare \emph{pattern} ricorrenti\footnote{C. M. Bishop, \emph{Pattern Recognition and Machine Learning}, Springer, 2006.} all'interno di grandi quantità di dati, replicando la logica dello \emph{stare decisis}. La distinzione tra algoritmo e intelligenza artificiale\footnote{\emph{Ex multis}, vedi i seguenti contributi sull'avvento dell'I.A. dal punto di vista filosofico, etico, tecnologico, giuridico: J.C.R. Licklider, \emph{Man-computer symbiosis,} in \emph{IRE Transactions on Human Factors in Electronics,} 1960, n.~1, pp.~4-11; R. Kurzweil,\emph{The Age of Spiritual Machines: When Computers Exceed Human Intelligence}, New York, 1999; N. Bostrom, \emph{Ethical Issues in Advanced Artificial Intelligence}, Cambridge University Press, 2003; D. Gunning, D. W. Aha, \emph{DARPA'sExplainableArtificial Intelligence (XAI) Program}, in \emph{Artificial Intelligence Magazine}, 2019, pp.~44 - 58; A. Agnihotri - S. Bhattacharya, \emph{Neuralink. Invasive neurotechnology for human welfare}, Sage publications, Los Angeles, 2023; F. J. Santamaría Ramos, \emph{Implicacioneséticas y legalesrelacionadas con el uso creciente de la IA}, in \emph{L'Ircocervo,} 23, n.~1/2024, pp.~139-174, T. García Berrio, \emph{Ética de la virtud para una IA compasiva}, ivi, pp.~318-340, et alii.}, valorizzata sia dal Regolamento europeo 2024/1689sull'IA\footnote{Per un'analisi del \emph{Digital Services Act} e dell'\emph{Artificial Intelligence Act}, v. G. Contissa, F. Galli, \emph{La governance della ``disinformazione aumentata'' tra Digital Services Act e AI Act}, in \emph{Federalismi.it}, n.~15/2025, ove gli Autori esaminano l'integrazione dell'intelligenza artificiale nei processi informativi e la conseguente emersione del fenomeno della ``disinformazione aumentata'', caratterizzato dalla produzione e diffusione su larga scala di contenuti falsi mediante sistemi automatizzati, nonché dall'impiego delle medesime tecnologie per finalità di individuazione e contenimento del fenomeno.} sia dalla legge nazionale n.~132 del 23 settembre 2025, è fondamentale: l'algoritmo è una sequenza deterministica di istruzioni, mentre l'intelligenza artificiale è un sistema capace di apprendere dati tramite il \emph{machine learning} e generare output non rigidamente predeterminati\footnote{Sull'introduzione di tecniche di \emph{machine learning} in plurimi settori, vedi I. Goodfellow, Y. Bengio, A. Courville, \emph{Deep Learning}, MIT Press, 2016.}. Nello specifico, l'A.I. act parla di ``un sistema basato su una macchina che è progettata per operare con diversi livelli di autonomia e che può generare output come previsioni, raccomandazioni o decisioni che influenzano ambienti fisici o virtuali'\,'\footnote{Regolamento UE 2024/1689.}. L'intelligenza artificiale può quindi definirsi come la capacità di un sistema tecnologico di fornire prestazioni assimilabili a quelle dell'intelligenza umana: pensare e agire secondo i modelli tipici della mente e del comportamento della creatura pensante\footnote{C. Piccinni, L. Bianco,\emph{L'intelligenza artificiale in prospettiva giuridica: una risorsa o un nocumento per la giustizia?}, in \emph{Quotidiano legale}, 2024, §2.; A. Tsamados, L. Floridi, M. Taddeo, \emph{Human control of AI systems: from supervision to teaming}, in \emph{AI and Ethics}, 2024, pp.~1-14.}. I sistemi di intelligenza artificiale, si differenziano a seconda dell'algoritmo, o programma, utilizzato. Va tenuto a mente che l'algoritmo può essere adoperato per costituire un'IA, e che un'IA può comporre algoritmi, ma non sono termini intercambiabili. Algoritmo non è sinonimo di intelligenza artificiale. Tuttavia, nell'ampia varietà degli algoritmi, alcuni sono usati dall'essere umano per programmare un software IA\footnote{Per esigenze di semplificazione definitoria, una nozione di intelligenza artificiale di frequente richiamata nella prassi è quella contenuta nell'Executive Order del Presidente degli Stati Uniti del 30 ottobre 2023 (sec.~3, lett. b). In tale sede, l'Amministrazione Biden definisce l'IA come ``\emph{a machine-based system that can, for a given set of human-defined objectives, make predictions, recommendations, or decisions influencing real or virtual environments}''.}. In primo luogo, l'algoritmo deterministico opera secondo una logica causalistica, in forza della quale a un determinato dato in ingresso corrisponde un unico passaggio successivo e, quindi, una sola risposta: si tratta di una scelta netta, ragionevolmente certa, priva di pluralismo decisionale. Diversamente, l'algoritmo probabilistico ammette una pluralità di passaggi successivi e restituisce più risposte possibili, rivelandosi particolarmente utile sotto il profilo strategico e maggiormente idoneo a interpretare situazioni complesse caratterizzate da numerose variabili. Infine, l'algoritmo generativo si fonda su una logica probabilistico-adattativa e mira a individuare la soluzione migliore possibile attraverso l'analisi e la rielaborazione di una vasta mole di dati in continuo aggiornamento, mediante tecniche di apprendimento progressivo. È quest'ultima tipologia che consente all'intelligenza artificiale di esprimere le prestazioni più avanzate, grazie al ricorso al \emph{machine learning}\footnote{N. Durante, \emph{La discrezionalità amministrativa ed i vizi del procedimento amministrativo nell'epoca dell'intelligenza artificiale}, in \emph{\href{http://www.giustizia-amministrativa.it}{www.giustizia-amministrativa.it},} 2024, p.~1.}. La ricostruzione, seppur necessariamente sintetica, introduce al tema dirimente della spiegabilità dell'intelligenza artificiale, ossia l'interpretabilità delle risposte fornite alla luce dei criteri adottati. Il concetto in sé è di difficile definizione, ma se ne segnalano alcune illustrazioni efficaci ed autorevoli: \emph{''Interpretabilityis the degree to which a human can understand the cause of a decision.''}\footnote{T. Miller, \emph{Explanation in Artificial Intelligence: Insights from the Social Sciences}, in \emph{Artificial Intelligence}, 2019, pp.~1 ss.}; \emph{''Interpretability}is \emph{the degree to which a human can consistentlypredict the model\textquotesingle sresult."}\footnote{B. Kim, R. Khanna, O. O. Koyejo, \emph{Examples are notenough, learn to criticize! Criticism for interpretability,} in \emph{NIPS\textquotesingle16: Proceedings of the 30th International Conference on Neural Information Processing Systems,} 2016, p.~2288 ss.}\emph{.} In sintesi, maggiore è il grado di interpretabilità del modello di apprendimento, più agevole sarà comprendere perché una previsione è stata fatta, o una certa decisione è stata presa. Il software non dovrebbe limitarsi a fornire soluzioni, ma dovrebbe anche essere in grado di giustificare il percorso logico che ha condotto al risultato proposto\footnote{Sul tema dell'ingegno dell'IA riconosciuto implicitamente dalla l. n.~132/2025 attraverso una modifica della l. n.~632/1941, fondamentale il contributo di F. Romeo, G. Verde, \emph{La valutazione della prova finale alla luce dell'AI Act: cosa cambia e cosa cambierà,} in AA. VV. (a cura di), \emph{Accertamento Penale, AI, e cybersicurezza}, Napoli, 2026, p.~230: \emph{``\emph{I sistemi artificiali sono quindi dotati di ingegno, volendo essere pienamente normativisti. Questa presa d'atto del legislatore coincide con quanto gli orizzonti filosofici più rivolti verso gli studi sulla mente e quelli delle scienze cognitive, essenzialmente americani, stanno da tempo affermando: non ci troviamo più di fronte ad un solo tipo di intelligenza che calcola, valuta, giudica e decide. L'intelligenza umana è accompagnata, in tutte le sue espressioni e facoltà, da quella dei sistemi artificiali. Come per gli umani sono sempre possibili differenze o divergenze nelle fasi valutative o decisorie, così tra umani e non umani è possibile che si verifichino altrettante differenze, come anche tra non-umani e non-umani. Cambia però il riferimento attorno al quale costruire un eventuale accordo; non c'è più quella somiglianza che fa porre gli individui partecipanti al confronto ed all'accordo come tutti appartenenti alla medesima classe, quella degli umani. Qui le classi sono due, umani e non-umani}''.}}. Proprio per questo si afferma il settore di ricerca della XAI, quell'\emph{Artificial Intelligence} che mira ad essere \emph{eXplainable}. Se un tempo l'obiettivo era meramente produrre un'IA funzionante, e tanto bastava a celebrare l'ingegno antropico, oggi la letteratura scientifica si sta concentrando sui modi per ingenerare fiducia in questo nuovo strumento dell'ingegno umano. Tale fiducia si fonda proprio sulla spiegabilità dei ragionamenti complessi di cui il programma è capace, e sulla differenza ingente tra i modelli a scatola nera e quelli a scatola bianca\footnote{C. Molnar, \emph{Interpretable Machine Learning. A Guide for Making Black Box Models Explainable}, Munich, 2022, p.~11: \emph{``A Black Box Model is a system thatdoesnotrevealitsinternalmechanisms. In machine learning,''black box'' describes models that cannot be understood by lookingattheirparameters (e.g.~a neural network). The opposite of a black box is sometimesreferred to as White Box, and isreferred to in this book asinterpretable model. Model-agnosticmethods for interpretabilitytreat machine learning models as black boxes, evenifthey are not''}.}. La \emph{white box} è quella trasparente, interpretabile, comprensibile. Purtroppo, sebbene un ragionamento effettivamente sussista dietro ogni input fornito, la sua trasparenza non è garantita per automatismo. Come le reti neurali umane, anche quelle artificiali sono formate da uno strato di ingresso, nascosto, e uno strato di uscita. È conoscibile quanto inserito in ingresso, cioè la richiesta dell'utente umano, ed è conoscibile la risposta in uscita sull'interfaccia digitale, ma occorre indagare sulla conoscibilità di quel che avviene nello strato nascosto. Le reti neurali profonde, che più assomigliano a quelle umane, sono delle vere e proprie \emph{black box}\footnote{V. F. Pasquale, \emph{The Black Box Society. The Secret Algorithms That Control Money and Information},Cambridge, 2015.}: svelare il processo intellettivo dietro l'erogazione di quella determinata risposta è oggetto di studio da parte della moderna linguistica computazionale. Da questa riflessione sorge un interrogativo tutt'altro che peregrino: perché non accontentarsi della correttezza del risultato, rinunciando a ogni indagine sulla sua interpretabilità? La risposta fornita da una parte significativa della dottrina è che il sacrificio della spiegabilità è tanto meno tollerabile quanto più rilevanti sono gli interessi coinvolti. Come è stato efficacemente osservato,''\emph{the problem is that a single metric, such as classification accuracy, is an incomplete description of most real-world tasks}''\footnote{C. Molnar, \emph{Interpretable Machine Learning}, cit., p.~11.}. Tale considerazione assume particolare rilievo nel settore pubblico, ove l'attività amministrativa incide per definizione su interessi pubblici e privati di rilevanza giuridica. In questo contesto, la mera accuratezza statistica del risultato non può ritenersi sufficiente: occorre che il percorso logico seguito dal sistema sia conoscibile e verificabile, così da consentire il controllo della decisione e l'effettiva tutela delle situazioni giuridiche coinvolte. Non può, poi, essere trascurato il problema dei bias, ossia dei pregiudizi che i sistemi di intelligenza artificiale possono ereditare dai dati utilizzati per il loro addestramento. Poiché l'apprendimento automatico si fonda sull'analisi di informazioni prodotte dagli esseri umani, i modelli tendono inevitabilmente a riprodurne rappresentazioni, preferenze e distorsioni. Ne consegue che i bias possono derivare tanto dalle scelte operate nella progettazione del sistema quanto dalle caratteristiche dei dati impiegati per l'addestramento. Essi possono manifestarsi in forma esplicita, come nel caso di contenuti discriminatori, oppure in modo più sottile, attraverso la riproduzione di abitudini, pratiche e schemi sociali consolidati\footnote{Come osserva C. Molnar, \emph{op. cit}., p.~15, i modelli di \emph{machine learning} tendono ad assorbire i bias presenti nei dati di addestramento, con il rischio di produrre decisioni discriminatorie nei confronti di gruppi storicamente svantaggiati. In tale prospettiva, l'interpretabilità costituisce uno strumento essenziale per individuare e correggere tali distorsioni. L'Autore evidenzia, ad esempio, che un sistema automatizzato per la concessione del credito, pur ottimizzato a ridurre il rischio di insolvenza, potrebbe violare il divieto di discriminazione, poiché quest'ultimo rappresenta un ulteriore vincolo del problema decisionale non incorporato nella sola funzione di ottimizzazione.}. \emph{Paucis verbis,} si fa quindi prima a dire quando non occorra perseguire l'interpretabilità: quando il modello non ha impatto significativo di natura economica o sociale, e quando il problema è già ampiamente conosciuto. In queste due ipotesi, la mera accuratezza del modello è più che sufficiente a giustificare la sua non spiegabilità: finché offre soluzioni efficaci, non è necessario indagare come lo faccia. Non serve sottolineare come queste due eccezioni siano difficilmente riscontrabili nell'ambito dell'amministrazione della giustizia. Si è detto in dottrina che l'algoritmo presenta alcune caratteristiche fondamentali, tra cui si evidenziano il determinismo, la prevedibilità, la controllabilità e la neutralità funzionale. Esempi tipici di impiego di strumenti algoritmici includono i sistemi di calcolo automatico dei termini processuali, i software di gestione del ruolo e gli strumenti di estrazione e classificazione documentale. Si tratta di ambiti nei quali l'operatività è interamente riconducibile a procedure deterministiche e regole predefinite, senza margini significativi di apprendimento o adattamento. Proprio per tale ragione, in questi casi non rileva il tasso di ``successo'' o di accuratezza statistica del sistema, quanto piuttosto la correttezza della sua programmazione rispetto alle regole giuridiche sottostanti. Ne consegue che tali applicazioni appartengono più propriamente alla sfera dell'automazione algoritmica in senso stretto, piuttosto che a quella dell'intelligenza artificiale in senso proprio, la quale presuppone invece capacità di inferenza, generalizzazione e adattamento a contesti non completamente formalizzati. Di tali profili appare consapevole la giurisprudenza amministrativa. In particolare, può richiamarsi, a titolo esemplificativo, la pronuncia del TAR Lombardia -- Milano, Sez. II, 31 marzo 2021, n.~843, ove si rimarca la distinzione tra algoritmo e intelligenza artificiale. Secondo il tribunale, un algoritmo consiste in una sequenza finita di istruzioni ben definite, eseguibili in modo meccanico per ottenere un determinato risultato. L'intelligenza artificiale, invece, implica capacità di apprendimento, adattamento e interazione con l'ambiente, consentendo al sistema di modificare i propri comportamenti in funzione dell'esperienza acquisita. Ne deriva che la presenza di un algoritmo, per quanto automatico, non equivale alla presenza di IA, che richiede una forma di autonomia decisionale e valutativa non riducibile a semplici procedure predeterminate.

Le definizioni così richiamate assumeranno valore di riferimento concettuale nel prosieguo dell'analisi, al fine di evitare ripetizioni e garantire una costante uniformità terminologica tra ``algoritmo deterministico'' e ``intelligenza artificiale'' in senso proprio.

\section{L'utilizzo degli algoritmi nel procedimento amministrativo}\label{lutilizzo-degli-algoritmi-nel-procedimento-amministrativo}

L'impiego dell'intelligenza artificiale nel procedimento amministrativo trova oggi espresso fondamento nell'art. 14 della l. n.~132/2025, che ne legittima l'utilizzo da parte delle pubbliche amministrazioni nel rispetto dei principi di conoscibilità, tracciabilità e responsabilità della decisione.

In tale quadro, l'utilizzo di strumenti algoritmici appare, in linea di principio, compatibile con l'attività vincolata, mentre suscita maggiori criticità quando sia coinvolta l'attività discrezionale, caratterizzata da valutazioni non interamente predeterminabili\footnote{Con riferimento al rapporto tra intelligenza artificiale - e, più precisamente, tra algoritmo - e decisione amministrativa, lo stato dell'arte può essere ricostruito nei termini essenziali che seguono. Quando l'attività amministrativa si presenta come vincolata, l'impiego di sistemi algoritmici non sembra, in linea di principio, porre particolari criticità. L'algoritmo opera infatti secondo uno schema deterministico, nel quale ai presupposti rigidamente fissati dalla legge corrisponde un esito integralmente predeterminato dall'ordinamento. In tali ipotesi, l'automazione costituisce una mera trasposizione tecnica di un procedimento già definito sul piano normativo, sicché la decisione risulta riconducibile a regole previamente stabilite e verificabili. Tale conclusione non è tuttavia esente da rilievi critici. Perplessità in ordine a un utilizzo meramente acritico dell'intelligenza artificiale, anche nell'ambito dell'attività amministrativa vincolata, sono state espresse, tra gli altri, da R. Cavallo Perin e I. Alberti, \emph{Atti e procedimenti amministrativi digitali}, in \emph{Il Diritto dell'Amministrazione Pubblica digitale}, II ed., a cura di R. Cavallo Perin e D. U. Galetta, Torino, 2025, pp.~138-176. Gli Autori osservano come l'automazione decisionale non possa prescindere da adeguate garanzie di trasparenza, controllabilità e verificabilità del procedimento, pena il rischio di una sostanziale deresponsabilizzazione dell'azione amministrativa. Anche nei casi in cui l'esito sia normativamente predeterminato, permane infatti l'esigenza di assicurare la comprensibilità del percorso decisionale e la possibilità di sindacarne la correttezza. Assai più problematico appare, invece, il ricorso all'intelligenza artificiale nell'ambito dell'attività discrezionale. Su questo terreno, dottrina e giurisprudenza non hanno ancora raggiunto approdi pienamente condivisi. Un primo orientamento guarda con favore all'utilizzo di strumenti algoritmici anche in contesti caratterizzati da discrezionalità, valorizzandone la capacità di elaborare grandi quantità di dati, di incrementare la coerenza delle decisioni e di ridurre il rischio di arbitrarietà. In tale prospettiva, l'algoritmo assumerebbe una funzione di supporto all'attività amministrativa, senza tuttavia sostituire il decisore umano. Un diverso indirizzo manifesta, invece, maggiori cautele. Secondo tale impostazione, la discrezionalità amministrativa implica valutazioni qualitative, comparazioni di interessi e apprezzamenti assiologici che difficilmente possono essere tradotti in regole computabili. L'automatizzazione del processo decisionale rischierebbe pertanto di irrigidire l'esercizio del potere amministrativo, comprimendo quegli spazi di flessibilità e di adattamento al caso concreto che costituiscono l'essenza stessa della discrezionalità. Le persistenti incertezze interpretative derivano, in larga misura, sia dal diverso grado di fiducia riposto nelle capacità tecniche dei sistemi di intelligenza artificiale, specialmente quelli basati su tecniche di apprendimento automatico, sia dalla stessa indeterminatezza della nozione di intelligenza artificiale, frequentemente impiegata in accezioni eterogenee e non sempre sovrapponibili. Nondimeno, sembra oggi emergere una soluzione intermedia. Pur escludendosi una piena sostituibilità dell'elemento umano nelle decisioni discrezionali, si ammette il ricorso a strumenti algoritmici quali ausili tecnici in grado di supportare il processo decisionale senza esaurirlo. Resta in ogni caso imprescindibile la riconducibilità della decisione all'autorità amministrativa, cui compete la responsabilità finale della ponderazione degli interessi coinvolti e della scelta del risultato da perseguire. Per un'approfondita ricostruzione dei diversi orientamenti si rinvia ad A. Di Martino, \emph{Tecnica e potere nell'amministrazione per algoritmi}, Napoli, 2023, nonché a G. Clemente di San Luca e M. Paladino, \emph{Decisione amministrativa e intelligenza artificiale}, in \emph{Federalismi.it}, 4/2025, p.~106 ss. Nel più ampio quadro delle trasformazioni determinate dalla digitalizzazione nel diritto pubblico, un riferimento imprescindibile è rappresentato dall'opera di L. Torchia, \emph{Lo Stato digitale. Una introduzione}, II ed., Bologna, 2025, che offre una ricostruzione sistematica delle principali questioni connesse all'impiego delle tecnologie digitali nei processi decisionali e nell'organizzazione amministrativa. Nello stesso solco si colloca il contributo di B. Marchetti, \emph{Amministrazione digitale}, in B.G. Mattarella e M. Ramajoli (a cura di), \emph{Enciclopedia del diritto. I Tematici. Funzioni amministrative}, III, Milano, 2022, p.~75 ss., che inquadra il fenomeno della digitalizzazione amministrativa alla luce delle più recenti evoluzioni dell'azione pubblica.}.

Gli algoritmi e i sistemi di intelligenza artificiale devono tuttavia presentare adeguati requisiti di affidabilità e adeguatezza rispetto alla funzione perseguita.

Ne consegue che non ogni soluzione tecnologica è di per sé idonea a soddisfare le esigenze dell'azione amministrativa, imponendosi piuttosto una valutazione concreta delle diverse opzioni disponibili, tanto più in un contesto di rapida espansione e differenziazione del mercato dei sistemi di intelligenza artificiale\footnote{Tale esigenza risulta particolarmente evidente ove si consideri che gli strumenti algoritmici impiegati dalle amministrazioni sono frequentemente sviluppati da soggetti privati esterni all\textquotesingle organizzazione pubblica. Come evidenziato da B. Ponti, \emph{Il fornitore dell\textquotesingle algoritmo quale soggetto estraneo all'amministrazione}, in \emph{Rivista italiana di informatica e diritto}, 2024, n.~2, pp.~517-526, l'esternalizzazione della progettazione e della gestione degli algoritmi pone delicate questioni in termini di trasparenza, imputazione delle responsabilità e sindacabilità delle decisioni. Il rischio è infatti che una quota significativa del processo decisionale venga, di fatto, trasferita a soggetti estranei all\textquotesingle amministrazione, con possibili ricadute sulla piena conoscibilità delle logiche sottese alla decisione pubblica e sulla stessa effettività del controllo amministrativo e giurisdizionale. La centralità di tali problematiche emerge con evidenza anche alla luce della straordinaria diffusione delle tecnologie di intelligenza artificiale, fenomeno ormai rilevato in modo sistematico tanto dagli analisti OSINT quanto dai principali osservatori accademici. Tra i contributi più significativi si segnalano i report dell\textquotesingle Osservatorio Digital Innovation della School of Management del Politecnico di Milano, che offrono una lettura articolata e progressiva del fenomeno. In particolare, meritano menzione \emph{Artificial Intelligence nel 2025: mercato, adozione e trasformazione delle aziende} (5 febbraio 2026), \emph{La nuova architettura dei pagamenti tra interoperabilità, moneta digitale e intelligenza artificiale} (12 marzo 2026), \emph{La Generative AI in azienda: evidenze dal campo} (30 luglio 2025), \emph{L\textquotesingle Intelligenza Artificiale nel quotidiano, tra vita lavorativa e privata} (6 giugno 2025) e \emph{L\textquotesingle adozione dell\textquotesingle Artificial Intelligence nel 2024: Italia a confronto con l\textquotesingle Europa} (6 febbraio 2025). Considerati nel loro insieme, tali studi descrivono il passaggio dell\textquotesingle intelligenza artificiale da tecnologia sperimentale a infrastruttura pervasiva dei processi economici, organizzativi e decisionali, confermando la necessità di un approccio giuridico che non si limiti a valutare astrattamente l\textquotesingle ammissibilità dell\textquotesingle impiego degli algoritmi, ma che si interroghi concretamente sulle caratteristiche dei sistemi adottati, sui soggetti che li sviluppano e sulle garanzie che essi sono in grado di offrire.}.

Assodata la fisiologica richiesta di specifici requisiti, ci si domanda quali debbano concretamente essere, quantomeno in via generale. Un primo cardine riconosciuto dalla giurisprudenza maggioritaria è senza dubbio quello della conoscibilità del meccanismo di funzionamento: il suo ruolo fondamentale è così valorizzato da poter superare ogni resistenza motivata dalla segretezza commerciale. Di questo avviso pare essere il Consiglio di Stato, che richiede, ai fini della legittimità della decisione assunta con l'ausilio di tale strumento, che ``il meccanismo attraverso il quale si concretizza la decisione robotizzata (ovvero l'algoritmo) debba essere conoscibile, secondo una declinazione rafforzata del principio di trasparenza, che implica anche quello della piena conoscibilità di una regola espressa in un linguaggio differente da quello giuridico''\footnote{Cons. St., sez. VI, 13 dicembre 2019, n.~8472.}. E giungendo precisamente al bilanciamento tra esigenza di conoscibilità e segretezza del software, la medesima sentenza prosegue: ``non può assumere rilievo l'invocata riservatezza delle imprese produttrici dei meccanismi informatici utilizzati i quali, ponendo al servizio del potere autoritativo tali strumenti, all'evidenza ne accettano le relative conseguenze in termini di necessaria trasparenza''\footnote{Ibidem. La trasparenza si presenta, infatti, anche secondo la dottrina, come la risorsa indispensabile per governare la sfida tecnologica nei suoi approdi ultimi, che portano alla piena automazione decisionale: in tal senso, A. Tronci, \emph{Le nuove frontiere della trasparenza nell'amministrazione algoritmica}, in \emph{\href{http://federalismi.it}{\ul{Federalismi.it}},} n.~9/2025, p.~141 ss. Più in generale, sulla legalità algoritmica e sulle garanzie di conoscibilità, spiegabilità e sindacabilità delle decisioni automatizzate, v. E. Carloni, \emph{I principi della legalità algoritmica. Le decisioni automatizzate di fronte al giudice amministrativo}, in \emph{Diritto amministrativo}, 2020, pp.~273-304; Id., \emph{Critica dell\textquotesingle amministrazione artificiale}, Bologna, 2026.}. Sicché da un lato la trasparenza deve garantire una migliore \emph{governance} dell'algoritmo, dall'altro deve, per questo, mutare ed adattarsi al funzionamento algoritmico, connotato da una patina sottile ma concreta di inconoscibilità. La trasparenza costituisce, pertanto, un obiettivo di non agevole perseguimento, soprattutto alla luce di processi automatizzati che tendono a divenire sempre più opachi, come evidenziato nel paragrafo precedente con riferimento al fenomeno della \emph{black box}. I processi decisionali attuati mediante soluzioni di \emph{machine learning} e di \emph{deep learning} pongono infatti numerosi problemi quanto alla capacità di comprenderne modi e ragioni, anche tenuto conto degli input ricevuti. I sistemi di IA risentono di una opacità che dipende da alcune caratteristiche proprie del fenomeno\footnote{Si pensi ai dati, che sono processati con modalità che rendono non ripercorribile il processo decisionale; e si pensi al \emph{machine deep learning}, che risulta sottratto a logiche deterministiche e prevedibilità a priori ma anche, nelle sue forme più evolute, a quelle di riverificabilità a posteriori.}. Per affrontare il tema dell'impiego dell'intelligenza artificiale da parte della pubblica amministrazione, assumendo come dato acquisito il ricorso già diffuso a forme più semplici di automazione algoritmica, è necessario interrogarsi anche sullo statuto epistemologico di tali strumenti, persino nell'ipotesi in cui essi siano pienamente trasparenti e conoscibili. Tale profilo costituisce il criterio dirimente per stabilire se l'intelligenza artificiale sia effettivamente in grado di comprendere il dominio semantico entro il quale opera oppure se, al contrario, si limiti a produrre elaborazioni linguistiche sofisticate fondate su correlazioni e inferenze di natura meramente probabilistica. Illustri autori nazionali\footnote{Importanti contributi sul tema: J. Nudo, M. E. Pandolfo, E. Loru, M. Samory, M. Cinelli, W. Quattrociocchi, \emph{Generative exaggeration in LLM social agents: Consistency, bias, and toxicity, Online Social Networks and Media}, 51, (100344), 2026; W. Quattrociocchi, V. Capraro, M. Perc, \emph{Epistemological Fault Lines Between Human and Artificial Intelligence,} in \emph{arXiv}, 2025.} hanno sollevato la questione, partendo da un aneddoto filosofico, da cui deriva il termine ``epistemia''\footnote{Nella Grecia antica, il rapporto tra Socrate e i sofisti è tradizionalmente descritto come antagonistico. Il filosofo era solito accusarli di essere meri retori, capaci di sostenere qualsiasi tesi e di adattare il discorso a ogni posizione, quasi ``tuttologi'' ante litteram. Essi non sarebbero stati dotati dell'ἐπιστήμη (A. Nehamas, \emph{Episteme and Logos in Plato\textquotesingle s Later Thought}, \emph{Archiv für Geschichte der Philosophie}, vol.~66, n.~1, 1984, pp.~11--36), intesa come conoscenza certa e incontrovertibile, nettamente distinta dall'opinione e dall'esperienza pratica, come emerge nel celebre dialogo platonico della \emph{Repubblica} e, successivamente, nel pensiero aristotelico (Aristotele, \emph{De Anima} II, 5, 417b).}. Ecco che l'epistemia è correttamente definita \emph{``the illusion of knowledge that emerges when surface plausibility replaces verification. Indeed, delegating judgment to such systems mayaffect the heuristics underlyingevaluativeprocesses, suggesting a shift from normative reasoningtoward pattern-based approximation and raising open questions about the role of LLMs in evacuative processes''}\footnote{Come evidenziato da E. Loru, J. Nudo, N. Di Marco, A. Santirocchi, R. Atzeni, M. Cinelli, V. Cestari, C. Rossi-Arnaud e W. Quattrociocchi nello studio \emph{The simulation of judgment in LLMs}, i \emph{large language models} sono sempre più impiegati in processi di classificazione, valutazione e supporto decisionale, sino a incidere su decisioni istituzionali ad alto impatto. Tale evoluzione segna il passaggio dalla mera automazione di compiti alla delega di vere e proprie funzioni valutative a sistemi sociotecnici. Gli Autori sottolineano, tuttavia, i rilevanti rischi connessi a tale fenomeno: modelli addestrati su dati storicamente distorti possono infatti replicare e amplificare disuguaglianze sociali, specialmente in ambiti sensibili come il mercato del lavoro.}\emph{.} Trattasi quindi di una conoscenza apparentemente scientifica, ma concretamente solo linguistica. È quel che distingue, in sintesi, il perito tecnico dall'esperto oratore: il primo detiene il sapere scientifico del suo dominio, come può essere una materia giuridica; il secondo detiene il potere di formulare frasi di senso compiuto in qualsiasi campo, ad esempio fornendo risposte credibili a quesiti giuridici. Dibattere sull'epistemia significa quindi una sola cosa: ritenere che l'IA abbia, o meno, una conoscenza scientifica dei saperi richiesti\footnote{Un esempio pratico pare qui essere utile. Ci si chiede, cioè, se la macchina ci dica che il cielo è azzurro perché essa sa che è azzurro, o perché attraverso un calcolo statistico ha ritenuto probabile che ``azzurro'' sia la risposta che l'utente umano si aspetta. D'altronde, l'essere umano può affermare non solo il colore corretto del cielo, ma anche i precisi motivi per cui esso è tale; l'IA può solo affermare il colore che probabilmente il cielo dovrebbe avere, non godendo di occhi per poterlo percepire sensorialmente.}. L'elevata accuratezza del modello linguistico è tra l'altro un'arma a doppio taglio: da un lato simboleggia l'elevata qualità dello strumento, ma dall'altro rende agevole cedere alla tentazione di non indagare sull'episteme dietro ogni risposta fornita. Innanzi ad una formulazione retorica forbita e convincente, l'utente è naturalmente scoraggiato dall'indagare la veridicità sostanziale delle affermazioni. Si potrebbe quindi sostenere che il rischio risieda, a lungo andare, nel credere che il cielo sia azzurro perché l'IA ha fornito una risposta ben motivata, e non perché \emph{ictu oculi} esso sia concretamente azzurro nel momento presente\footnote{In tal senso, Parasuraman e Riley hanno evidenziato come l'interazione uomo-macchina possa indurre fenomeni di eccessiva dipendenza dall'automazione, con conseguente indebolimento del controllo umano effettivo (R. Parasuraman -- V. Riley, \emph{Humans and Automation: Use, Misuse, Disuse, Abuse}, in \emph{Human Factors}, vol.~39, n.~2, 1997, pp.~230--253). Più recentemente, il problema è stato ripreso con riferimento ai sistemi di intelligenza artificiale generativa, in relazione ai quali si è sottolineata la loro intrinseca opacità epistemica (\emph{black box problem}) e la difficoltà di ricostruire le ragioni sottese alle singole risposte prodotte (F. Pasquale, \emph{The Black Box Society: The Secret Algorithms That Control Money and Information}, Cambridge, MA, Harvard University Press, 2015).}. Dalla teoria che si vuol sposare derivano conseguenze inevitabili: se si ritiene che i modelli linguistici siano caratterizzati da una vera conoscenza epistemologica, allora essi sono certamente adoperabili nel settore pubblico, e a costoro è delegabile l'emanazione di provvedimenti vincolati e discrezionali; se si sposa la teoria dell'epistemia, allora qualsiasi modello è solo un discreto oratore capace di simulare la conoscenza di qualsiasi argomento, ed è dunque inadatto a qualsiasi funzione sensibile\footnote{Questo meccanismo opererebbe allo stesso modo verso un soggetto umano: il responsabile può emanare il provvedimento in quanto conosce scientificamente la sua essenza, la sua funzione, il suo contenuto, e non perché ricorda semplicemente a memoria una formula convincente di parole in sequenza.}. La mera apparenza non può aver posto nella funzione pubblica, e tantomeno nell'esercizio della discrezionalità amministrativa. Invero, il suo esercizio presuppone intanto un decisore in grado di valutare gli opposti interessi in campo, e non di valutare la loro astratta rappresentazione linguistica. Per tali ragioni, se l'IA è mera epistemia di massa, allora essa è totalmente inadeguata all'esercizio di qualsiasi discrezionalità, in quanto slegata da ogni reale comprensione della realtà umana. Alla luce di quanto sopra asserito, si ritiene che le intelligenze artificiali attualmente esistenti siano caratterizzate da una dimensione meramente conoscitiva e che, dunque, siano certamente artificiali, ma per nulla intelligenti. Se è palese che l'utilizzo delle intelligenze artificiali da parte dell'amministrazione possa avvenire solo con una maggiore alfabetizzazione digitale e cognitiva degli operatori, occorre domandarsi se tale utilizzo possa costituire un mero ausilio, come dai più asserito, o una completa sostituzione del funzionario antropomorfo, come sostengono in pochi\footnote{Non necessita di giustificazione la maggiore tolleranza verso la prima ipotesi, dato che uno strumento con elevato grado di accuratezza è \emph{in re ipsa} utile come ausilio tecnologico a chiunque lo adoperi. Altrettanto scontata appare la maggiore prudenza manifestata verso la seconda ipotesi: eliminare la presenza del decisore umano è già difficilmente accettabile nell'odierno contesto socioculturale, ma in particolar modo lo è quando il nuovo decisore artificiale potrebbe finanche non essere dotato di alcuna conoscenza scientifica del dominio di riferimento.}. Occorre, a questo punto, prospettare due ipotesi contrapposte, anche al fine di valorizzare i recenti apporti giurisprudenziali. Se i modelli linguistici sono dotati di conoscenza epistemologica, allora possono condividersi le affermazioni del Consiglio di Stato circa l'ammissibilità dell'esercizio del potere discrezionale mediante l'impiego dell'intelligenza artificiale; se, invece, tali modelli ne sono privi, il loro utilizzo dovrà essere limitato all'esercizio dei poteri vincolati. In questi termini, appare plausibile rinvenire il fondamento del favore manifestato dal Supremo organo della giustizia amministrativa nei confronti di algoritmi e modelli linguistici anche nell'adozione di provvedimenti discrezionali, purché essi svolgano una funzione di mero ausilio ai funzionari umani, i quali restano decisori finali e infungibili. In ordine al grado di aleatorietà, gli algoritmi più sicuri e prevedibili sono quelli deterministici, detti anche condizionali. Essi sono ampiamente impiegati sin dalla seconda metà del Novecento e risultano certamente ammissibili quali strumenti di supporto informatico della pubblica amministrazione. Tali algoritmi possono assumere decisioni, elaborare piani o formulare inferenze sulla base dell'applicazione di regole predeterminate nella fase di programmazione dell'algoritmo stesso\footnote{Vedi ad esempio l'algoritmo intelligente condizionale sui Criteri Ambientali Minimi relativi alla ``acquisizione di sorgenti luminose per illuminazione pubblica, l'acquisizione di apparecchi per illuminazione pubblica, l'affidamento del servizio di progettazione di impianti per illuminazione pubblica'', approvato con DM 27 settembre 2017. Ivi si prevede come criterio preferenziale il fatto che ``nella regolazione del flusso luminoso si tiene conto del valore di luminanza reale della strada illuminata, sia tramite misura della luminanza in tempo reale sia tramite algoritmi di presunto decadimento del flusso luminoso'' (p.~4.3.4.5). L'algoritmo condizionale entra così in azione: se la luminosità cala al di sotto di un certo valore di luminanza, allora l'illuminazione stradale deve accendersi, e non occorre che sia un essere umano a farlo, né prefissare un orario di accensione determinato.}. In altri termini, l'algoritmo esegue comandi che sono stati previamente stabiliti e che costituiscono essi stessi il modello procedurale del programma.\footnote{G. Carullo, \emph{Decisione amministrativa e intelligenza artificiale}, in \emph{Diritto dell\textquotesingle informazione e dell\textquotesingle informatica, fasc.~3, 2021, p.~434.}} All'inserimento delle regole informatiche provvedono di norma gli esseri umani, o finanche altri \emph{software} nell'ipotesi di IA programmate da IA. Evidentemente l'algoritmo condizionale ben si presta all'esercizio del potere vincolato, a patto che l'amministrazione sia capace di tradurre norme giuridiche in regole informatiche, e cioè in condizioni di innesco della macchina\footnote{Si pensi così all'erogazione di un beneficio non appena il destinatario sia divenuto idoneo a riceverlo, o all'accensione dell'illuminazione urbana quando la pubblica via sia divenuta eccessivamente oscura. In queste ipotesi la prima condizione sarà il raggiungimento del requisito per godere del beneficio, e la seconda un certo grado minimo di luminanza.}. Offrendo legittimità a tale metodo di realizzazione del pubblico agire, il Consiglio di Stato ha postulato che ``la regola tecnica che governa ciascun algoritmo resta pur sempre una regola amministrativa generale, costruita dall'uomo e non dalla macchina, per essere poi (solo) applicata da quest'ultima, anche se ciò avviene in via esclusiva''\footnote{Cons. St., Sez. VI, sent. n.~2270/2019.}. Se non vi sono dubbi circa l'ammissibilità degli algoritmi condizionali nell'esercizio dei poteri pubblici, occorre rilevare una significativa obiezione proveniente dall'ingegneria informatica: tale soluzione diviene impraticabile quando le condizioni da considerare sono troppo numerose o difficilmente quantificabili. L'esempio classico è quello dell'identificazione di un oggetto in un'immagine. Attraverso un algoritmo deterministico, sarebbe necessario inserire un elevatissimo numero di regole affinché il sistema riconosca l'oggetto nelle sue possibili dimensioni, varianti, forme e colori. Questa osservazione risulta particolarmente utile nel diritto amministrativo, in ragione dell'evidente incompatibilità tra algoritmo condizionale ed esercizio del potere discrezionale: occorrerebbe predisporre una quantità enorme di istruzioni per consentire all'IA di ponderare adeguatamente la decisione, con il limite ulteriore che il software così realizzato sarebbe utilizzabile soltanto per uno specifico potere amministrativo. In sintesi, si tratta del noto problema dei programmi dotati di \emph{hard-coded knowledge}, la cui conoscenza è limitata ai dati e alle regole espressamente inseriti dal programmatore\footnote{I. Goodfellow, Y. Bengio, A. Courville, \emph{Deep Learning}, MIT Press, Cambridge, Massachusetts, 2016, p.~2.}. Sufficienti per l'esercizio del potere vincolato, insufficienti per qualsiasi forma di potere discrezionale. L'incompletezza della programmazione condizionale, e dunque degli algoritmi deterministici, ha portato autorevoli voci della comunità scientifica ad asserire che\emph{``the difficulties faced by systems relying on hard-coded knowledge suggest that AI systems need the ability to acquiretheirown knowledge, by extracting patterns from raw data''}\footnote{Ibidem.}. Questa ambiziosa idea di sistema che si nutra autonomamente di informazioni da fonti proprie si è realizzata con l'avvento del \emph{machine learning}, illustrato in questi termini: \emph{``searching for useful representations and rules over some input data, within a predefinedspace of possibilities, usingguidance from a feedback signal''}\footnote{F. Chollet, \emph{Deep learning with Python}, II Ed., Manning, Shelter Island, New York, 2021, p.~2.}\emph{.} Suole dirsi che nel \emph{machine learning} l'obiettivo della macchina nella fase di apprendimento è riscontrare rappresentazioni di un certo set di dati conosciuto, ossia i \emph{training data}, sulla base delle quali successivamente analizzare dati ignoti\footnote{Sulla rilevanza dei dati, anche nel settore pubblico, vedi G. Carullo, \emph{Gestione, fruizione e}

  \emph{diffusione dei dati dell'amministrazione digitale e funzione amministrativa,} Giappichelli, Torino, 2017, p.~243; G. Carullo, \emph{Big Data e pubblica amministrazione nell'era delle banche dati interconnesse}, in Conc. Merc., vol.~23, 2016.} e quindi produrre un risultato predittivo attendibile\footnote{Durante l'apprendimento si forma il modello, il concetto base, l'idea \emph{paucisverbis}, ed esso viene utilizzato \emph{in executivis}dell'algoritmo per verificare se un set di nuovi dati si conformi ad esso. Il quesito algoritmico è se i nuovi dati possano essere ricondotti al medesimo concetto, cioè se non presentino peculiarità originali che li rendano estranei al modello.}. La possibilità di scegliere tra più opzioni ugualmente legittime e il carattere aleatorio di tale scelta rendono gli algoritmi non \emph{hard coded} astrattamente più adatti all'impiego nell'esercizio dei poteri discrezionali. Questi ultimi si esercitano in quei campi nei quali non esiste un'unica soluzione corretta, né una sola via per il perseguimento di un obiettivo, e in cui le decisioni non conseguono al mero avverarsi di una determinata condizione. In questa prospettiva, il tasso di attendibilità assurge a parametro di misurazione dell'idoneità della risposta fornita a risolvere il problema sottoposto alla macchina. Poiché il processo di analisi si fonda su rappresentazioni matematico-numeriche, la deviazione del risultato rispetto al modello generato nella fase di apprendimento può essere quantificata e sottoposta a limitazioni. La pubblica amministrazione potrebbe così arginare gli errori nell'esercizio del potere discrezionale, imponendosi di adottare soluzioni generate algoritmicamente soltanto qualora siano caratterizzate da un elevato tasso di attendibilità, ossia da una minima percentuale di errore e da un elevato grado di accuratezza\footnote{Un esempio può rinvenirsi nella pianificazione urbanistica o nella localizzazione di un'infrastruttura pubblica. In tali ipotesi, la pubblica amministrazione è chiamata a comparare una pluralità di interessi pubblici e privati tra loro concorrenti (tutela ambientale, sviluppo economico, mobilità, impatto sul territorio), senza che l'ordinamento predetermini in modo puntuale un'unica soluzione obbligata. In questo contesto, un sistema di IA potrebbe elaborare differenti scenari decisionali, simulando gli effetti delle diverse opzioni possibili e restituendo alla pubblica amministrazione quella che, secondo i parametri modellizzati, presenta il più elevato tasso di attendibilità rispetto agli obiettivi pubblici perseguiti, ferma restando la decisione finale in capo al soggetto umano.}. In tale prospettiva, l'IA può essere impiegata anche nell'ambito dei procedimenti discrezionali, ma esclusivamente quale strumento di supporto all'attività decisionale. Laddove il contributo umano non sia sacrificabile, l'IA non può sostituire i pubblici impiegati, ma può comunque offrire proposte di soluzione dei problemi in tempi celeri e tendenzialmente affidabili. La discrezionalità in quanto tale resta prerogativa dell'essere umano, ma ciò non esclude la possibilità di prestare orecchio ad un'\emph{IA-based advisor}, un erogatore di consigli e suggestioni di natura informatica. Emerge così chiaramente la differente utilità di algoritmi condizionali, a cui è delegabile\footnote{Lo stesso termine ``delegare'' risulta, tuttavia, improprio: la pubblica amministrazione che utilizza un algoritmo non trasferisce ad esso alcun potere, in quanto questo è attribuito dalla legge all'organo competente e non agli strumenti di cui esso si avvale. Ciò rileva sul piano dell'imputabilità delle decisioni algoritmiche, nonché in relazione al rischio che il funzionario pubblico si limiti a seguire pedissequamente le indicazioni della macchina, adottando una sorta di ``pilota automatico''. Una simile impostazione si tradurrebbe, di fatto, in una sostituzione integrale del funzionario umano con un decisore automatizzato, ipotesi rispetto alla quale ci si è già espressi in senso negativo.} l'esercizio del potere vincolato, e di algoritmi di \emph{machine learning}, funzionali, nei limiti sopra individuati, all'esercizio del potere discrezionale. Il Consiglio di Stato non ha escluso \emph{``\emph{la riferibilità della decisione finale all'autorità ed all'organo competente in base alla legge attributiva del potere}''}\footnote{Cons. Stato, Sez. VI, 13 dicembre 2019, sent. n.~8472, 8473 e 8474.} in un caso riguardante gli algoritmi condizionali adoperati in potere vincolato dalla p.a., e ha ritenuto che eventuali errori valutativi sarebbero ugualmente contrastabili con la tutela giurisdizionale come qualsiasi errore umano. Il problema più pressante sembrerebbe invece essere quello della corretta formulazione delle regole informatiche applicate dall'algoritmo, e quindi la possibilità di verificare il modello procedurale seguito dal \emph{software} per giungere ad una determinata decisione\footnote{Il che, nella prospettiva del diritto amministrativo, si traduce in una questione attinente alla correttezza dell'iter logico-giuridico che ha portato alla formazione dell'atto e quindi alla relativa motivazione.}. Diversa è la questione dell'imputabilità delle decisioni assunte mediante algoritmi di machine learning\footnote{Mentre negli algoritmi deterministici le regole informatiche sono scritte da un agente sulla base delle norme giuridiche da applicare alla fattispecie, qui non vi è una codificazione delle regole che andranno a determinare le decisioni della macchina.}. In tali sistemi non operano regole predeterminate e rigidamente codificate, bensì modelli matematico-statistici elaborati a partire dai dati di addestramento, finalizzati a individuare correlazioni e schemi utili all'identificazione o alla classificazione di determinate entità. Il funzionamento del \emph{machine learning} si fonda, infatti, sull'apprendimento di un concetto attraverso l'analisi di una pluralità di esempi: una volta acquisito il modello rappresentativo di una determinata categoria -- ad esempio quella di ``mela'' -- il sistema è in grado di riconoscere una specifica mela, distinguendola sia da altri frutti sia dagli elementi presenti nello sfondo dell'immagine. Ne consegue che il sistema genera autonomamente i parametri destinati a governare la fase esecutiva, costruendo un modello decisionale sulla base del patrimonio informativo utilizzato nell'addestramento\footnote{F. Maggino, G. Cicerchia, \emph{Algoritmi, etica e diritto}, in \emph{Dir. informaz. e inf}., 6, 2019, p.~1162, ben evidenzia che\emph{``\emph{il valore dei dati dipende molto dalla capacità di chi li analizza di estrarre informazioni e di ottenere conoscenza}''}.}. Pertanto, nel caso del machine learning, il problema dell'imputabilità non si concentra tanto sull'applicazione di una regola previamente definita, quanto piuttosto sulla costruzione del modello stesso, e dunque sulle scelte compiute nelle fasi di progettazione, addestramento, selezione e trattamento dei dati. È, infatti, in tale fase che vengono definiti, seppur in modo non completamente prevedibile e controllabile \emph{ex ante}, le correlazioni, i pesi e i criteri che orienteranno le successive elaborazioni del sistema. In questa prospettiva, la decisione finale resta comunque imputabile alla pubblica amministrazione, la quale non si limita a recepire passivamente un esito automatizzato, ma sceglie di avvalersi di uno strumento che produce inferenze sulla base di un patrimonio informativo preesistente. Ne consegue che la responsabilità dell'azione amministrativa si estende anche alle scelte relative alla progettazione del sistema, alle modalità di addestramento e alla selezione dei dati impiegati, poiché la qualità e la struttura di tali elementi incidono direttamente sull'output generato dall'algoritmo. Alla luce delle considerazioni sin qui svolte, sembra corretto affermare che non esista un divieto generale di ricorso all\textquotesingle intelligenza artificiale nel procedimento amministrativo e che, conseguentemente, il suo impiego non richieda una specifica autorizzazione contenuta nella singola norma attributiva del potere. Tale conclusione era già stata prospettata con riferimento agli algoritmi deterministici e trova fondamento in un orientamento dottrinale consolidato, formatosi ben prima dell'attuale sviluppo delle tecnologie di intelligenza artificiale\footnote{Si segnala a riguardo A. Masucci, \emph{L'atto amministrativo informatico: primi lineamenti di una ricostruzione}, Jovene, Napoli, 1993, pp.~53--54, il quale afferma che il \emph{``\emph{fondamento giuridico del programma}''} può essere individuato nella \emph{``\emph{potestà autorganizzatoria}''} della p.a.; D. U. Galetta, \emph{Algoritmi, procedimento amministrativo e garanzie}, cit., par. 2.}. La tesi che sostiene l'implicita legittimazione della p.a. a dotarsi di propria IA a fini istituzionali\footnote{G. Carullo, \emph{Decisione amministrativa e intelligenza artificiale}, cit. p.17 ss. e in generale l'intero par. 6, titolato ``\emph{L'implicita legittimazione all'uso di strumenti di ia nell'attribuzione del potere.''}} è forte di queste argomentazioni storicamente risalenti all'avvento dei primi algoritmi, oltre che sostenuta dai sopra richiamati precedenti della giustizia amministrativa\footnote{\emph{Ex multis}, v. Cons. St., Sez. VI, 4 febbraio 2020, n.~881; Id., 13 dicembre 2019, n.~8472, 8473 e 8474; Id., 8 aprile 2019, n.~2270.}, che negano l'incompatibilità delle intelligenze artificiali con le normali garanzie procedurali ordinariamente poste a tutela del cittadino\footnote{Anzi, sostengono alcuni che l'utilizzo dell'IA possa finanche incrementare il grado di trasparenza del processo decisionale: la mente umana del funzionario non è composta da righe intelligibili di codice, ma l'IA lo è. I bias cognitivi, i dati di allenamento, i c.d. pregiudizi e le loro cause sono dunque noti quando affliggono l'IA, mentre resteranno per sempre ignoti quando confinati nel sistema nervoso centrale del responsabile del procedimento. In questi termini, pare forse più sindacabile l'esercizio del potere discrezionale consigliato da una IA rispetto a quello compiuto dal nostro dipendente pubblico, oltre che dotato di maggiori garanzie di imparzialità.}. Ne consegue che l'amministrazione può legittimamente avvalersi - e già ampiamente si avvale - di sistemi algoritmici e di strumenti di intelligenza artificiale, quantomeno nelle attività istruttorie, organizzative e di supporto alla decisione. Diversamente, non appare allo stato compatibile con i principi fondamentali dell'ordinamento una totale sostituzione del decisore umano mediante processi integralmente automatizzati. Le peculiarità delle funzioni pubbliche, e in particolare delle attività che implicano valutazioni discrezionali o ponderazioni di interessi, impongono infatti il mantenimento di un presidio umano effettivo, al quale restino imputabili la decisione finale e la relativa responsabilità. In tale direzione si è mosso il legislatore con il citato art. 14 della l. n.~132/2025, disposizione che recepisce e sistematizza approdi già consolidati in dottrina e giurisprudenza, prevedendo che le pubbliche amministrazioni possano utilizzare sistemi di intelligenza artificiale al fine di migliorare l'efficienza, la qualità e la tempestività dell'azione amministrativa. La disciplina conferma, pertanto, l'adozione di un modello antropocentrico di amministrazione algoritmica. L'intelligenza artificiale è infatti configurata come strumento di supporto e di ausilio all'attività amministrativa, mentre l'autonomia decisionale e la responsabilità giuridica restano integralmente imputate alla persona fisica titolare del potere. Ne consegue che la legittimità dell'impiego di tali tecnologie non dipende dalla loro capacità di sostituire il funzionario, bensì dalla loro attitudine a coadiuvarne l'attività nel rispetto dei principi di legalità, trasparenza, buon andamento e responsabilità amministrativa.

\section{Intelligenza artificiale e funzione giurisdizionale}\label{intelligenza-artificiale-e-funzione-giurisdizionale}

Delineata l'utilità dello strumento algoritmico nell'ambito del procedimento amministrativo, si procede all'esame del successivo profilo: è auspicabile l'introduzione dell'intelligenza artificiale nella funzione giurisdizionale in generale? In termini sintetici, può ritenersi condivisibile un impiego dell'IA in funzione ausiliaria all'attività giurisdizionale, mentre deve escludersi la sua sostituzione al giudice, tanto nella giurisdizione ordinaria quanto in quella amministrativa. Tale conclusione trova fondamento, a livello interno, nei principi costituzionali del giusto processo (art. 111 Cost.) e della soggezione del giudice soltanto alla legge (art. 101 Cost.), nonché, a livello eurounitario, nel principio di controllo umano significativo sancito dall'art. 22 del Regolamento (UE) 2016/679 (GDPR) e nei principi di affidabilità, trasparenza e responsabilità richiamati dal Regolamento (UE) 2024/1689 (AI Act), che esclude l'utilizzo di sistemi di intelligenza artificiale in grado di sostituire integralmente la decisione giudiziaria umana. Il ragionamento compiuto segue la stessa trama tracciata nel paragrafo precedente, in cui sposando i recenti approdi legislativi, si vuol sostenere un'ampia implementazione degli strumenti algoritmici come ausilio del protagonista umano, che non può mai degradare ad attore secondario nell'esercizio di alcun potere pubblico, né giudiziario, né di polizia\footnote{Si segnala il longevo successo di \emph{Keycrime}, oggi \emph{Giove}, progetto nato dall'assistente di polizia Mauro Venturi nel 2008, e tutt'oggi in corso di sviluppo, che culminerà con la pubblicazione internazionale. Il suo successo risiede nel crime linking, e cioè non nella tradizionale analisi statistica degli hotspot a maggior incidenza delinquenziale, bensì nella profilazione dei pochi autori che commettono la maggioranza dei reati. Ivi il crime linking consiste nel collegare il crimine al criminale, e non le aree depresse alla criminalità. Il sistema non agisce in sostituzione degli apparati securitari, si limita ad agevolarne la distribuzione di risorse investigative. \href{https://www.ilsole24ore.com/art/intelligenza-artificiale-keycrime-giove-nuovo-investigatore-la-polizia-AE9i1EbD}{\ul{Intelligenza artificiale: da Keycrime a Giove, nuovo investigatore per la Polizia - Il Sole 24 ORE}}.}. Come l'IA non può compiere valutazioni in luogo del funzionario pubblico, così essa non può abbracciare o meno una teoria giuridica al posto del giudice.\footnote{Le questioni relative alla rilevanza della fase del processo suscettibile di automazione sono esaminate in modo approfondito da D. Barysė e R. Sarel, \emph{Algorithms in the Court: Does It Matter Which Part of the Judicial Decision-Making IsAutomated?,} in \emph{Artificial Intelligence and Law}, 32, 2024, pp.~117--146, nonché da G. Yalcin, E. Themeli, E. Stamhuis, S. Philipsen e S. Puntoni, \emph{Perceptions of Justice by Algorithms}, in \emph{Artificial Intelligence and Law}, 31, 2023, pp.~269--292.Da tali contributi emerge come, tra le possibili applicazioni dell'intelligenza artificiale in ambito giuridico, quelle volte a supportare il giudice nell'interpretazione delle norme, nell'argomentazione, nella costruzione del percorso logico e nell'adozione della decisione siano generalmente accolte con maggiore diffidenza rispetto agli strumenti destinati all'accertamento dei fatti o alla valutazione di elementi percettivi.}

\subsection{L'IA come strumento accessorio del giudice}\label{lia-come-strumento-accessorio-del-giudice}

In quanto mero meccanismo di ausilio, è infatti certamente auspicabile attribuire alcuni compiti materiali e logistici ai \emph{software} di \emph{machine learning,} tra cui l'organizzazione delle udienze, la gestione dei ruoli e la ricerca intelligente di precedenti giurisprudenziali. Tale impiego trova fondamento nei principi di buon andamento e organizzazione efficiente dell'amministrazione (art. 97 Cost.), nonché nel carattere personale della funzione giurisdizionale (artt. 101 e 111 Cost.), che esclude la delega decisoria ma non l'utilizzo di strumenti di supporto\footnote{A livello eurounitario, la compatibilità di tali impieghi è confermata dal Regolamento (UE) 2024/1689 (AI Act), che ammette sistemi di intelligenza artificiale a funzione ausiliaria, e dal principio di controllo umano significativo desumibile dall'art. 22 del GDPR.}. La attività sopra menzionate sono senza dubbio idonee ad essere svolte da sistemi algoritmici non umani, giacché non coperti da alcuna c.d. riserva di umanità ontologica: le decisioni assunte su questi aspetti si fondano su meri criteri utilitaristici di convenienza matematica, e certamente ancorati ad aspetti probabilistici. Si pensi alla fissazione delle udienze nel ruolo, ove ciascun membro del collegio segna manualmente date e orari. Non sussiste alcuna ragione giuridica, filosofica o scientifica per la quale un'IA non sia in grado di compilare le caselle cronologiche organizzative in via informatica. Similmente bisognerebbe anche dire che non vi è alcuna ragione per conservare un ruolo cartaceo, dato che la sua conservazione e compilazione è più efficiente nel formato digitale. Si consideri, inoltre, l'attività di disvelamento delle cause prive di particolare complessità, affidata all'Ufficio per il processo ai sensi dell'art. 72-\emph{bis}. In tale contesto, anche l'individuazione dei precedenti giurisprudenziali, funzionale all'adozione della sentenza in forma semplificata, può essere efficacemente demandata all'intelligenza artificiale. Lo stesso vale per la ricerca di sentenze, articoli e monografie. Così i precedenti rinvenuti dall'IA non sarebbero diversi dalle sentenze che un magistrato umano può oggi svelare attraverso una ricerca manuale su un \emph{browser,} dato che entrambi i risultati sono dovuti ad un riscontro di parole chiave inserite nel \emph{prompt} di ricerca. Ed anzi, la ricerca a mezzo IA quantomeno renderebbe uniforme e trasparente l'esito di un procedimento che, ad oggi, è compiuto in solitaria dal magistrato, e di cui sovente non resta traccia alcuna. Per ovviare al problema si potrebbe provvedere nel senso indicato da Simonetti\footnote{H. Simonetti, \emph{Brevi note}, cit., p.~929: \emph{``Per la dottrina, più che per le sentenze dei giudici, le sue opinioni non si contano ma si pesano e c'è una buona dose di soggettività nella loro selezione.Nella ricerca dei criteri sulla base dei quali scegliere quale dottrina ricercare nel vasto mare delle riviste, molte delle quali sono oramai solo on-line, sarebbe consigliabile fare affidamento sulle fasce elaborate all'interno del sistema universitario e magari limitarsi solo alla più alta di quelle fasce.''}}, limitando il campo d'azione soprattutto nella ricerca di dottrina, ove non conta tanto il numero di pubblicazioni riscontrabili, bensì la qualità di esse. Lo stesso si dica per l'applicazione di qualsiasi altro filtro di pertinenza alla ricerca, che può essere inserito dall'uomo e dalla macchina in egual misura. Il vero \emph{discrimen} risiede nella velocità di ricerca ed elaborazione dei dati: il magistrato impiegherebbe un tempo notevolmente maggiore rispetto a quello richiesto dall'IA per portare a compimento la medesima azione. E più meramente meccanica è l'azione, più sarebbe utile che fosse un software ad occuparsene, sgravando del carico i giudicanti, che possono concentrarsi sulle attività intellettivamente più impegnative. Si vuol pertanto sposare il prudente approdo a cui è giunto il legislatore con la citata legge n.~132/2025, in base alla quale\emph{``nei casi di impiego dei sistemi di intelligenza artificiale nell\textquotesingle attività giudiziaria è sempre riservata al magistrato ogni decisione sull\textquotesingle interpretazione e sull\textquotesingle applicazione della legge, sulla valutazione dei fatti e delle prove e sull\textquotesingle adozione dei provvedimenti.} La stessa legge rimarca inoltre la strumentalità dei sistemi di intelligenza artificiale rispetto all'organizzazione dei servizi relativi alla giustizia, ove spetta al Ministero della Giustizia procedere a semplificazioni\footnote{L. n.~132/2015, art. 15, co. 1-2.}.

\subsection{L'interprete automa}\label{linterprete-automa}

Differente e più complessa tematica è quella della sostituzione del magistrato con l'intelligenza artificiale nell'attività interpretativa ed innovativa, e cioé nell'ermeneutica giuridica e negli \emph{overruling}. Questo quesito ha certamente terreno fertile negli atteggiamenti maggiormente ``creativi'' dei giudici di merito, che non passano inosservati agli occhi dei giuristi. D'altronde, come osservato in dottrina ``\emph{una certa creatività giurisprudenziale può ritenersi legittima laddove intervenga a colmare lacune o anomie dell'ordinamento, garantendo coerenza e continuità al sistema giuridico. Diversamente, risulta difficilmente giustificabile quando il giudice, oltrepassando i confini della funzione giurisdizionale, sovrappone la propria valutazione discrezionale a quella già espressa dal legislatore, incidendo indebitamente sull'ambito riservato alla discrezionalità politica}''\footnote{G. Taglianetti, \emph{Sul confine tra creatività e creazionismo giurisprudenziale. considerazioni a margine della sentenza TAR Lombardia, sez. III, 18 aprile 2025, n.~1400, in tema di esame di abilitazione forense e obbligo di motivazione}, in \emph{Nomos. Le attualità del diritto}, n.~2/2025, p.~12.}. Ed è anche vero che la tendenza di ricercare un metodo per automatizzare le decisioni giudiziarie si è già manifestata nella tentata ideazione di un giudice-automa in Unione Sovietica nel lontano 1964\footnote{Esperimento della Sezione siberiana dell'Accademia delle Scienze dell'Unione Sovietica diretto ``alla creazione di un giudice automatico, cioè di un automa giudicante'', come attestato in V. Frosini, \emph{Informatica, diritto e società,} Milano, 1992, pp.~278-279.}, il che dimostra che la domanda non sia peregrina. La crisi di fiducia nel magistrato umano è strettamente correlata all'avvento dell'IA, giacchè l'ambizione di creare il giudicante robotico non può che derivare dall'insoddisfazione generata dal giudicante antropomorfo. Si condividono qui le parole di un autorevole membro della giurisdizione amministrativa, secondo cui si è affermato \emph{``\emph{un modello di giurisprudenza incline al soggettivismo e alla creatività: sempre più riluttante a fare applicazione delle semplici disposizioni -- che tanto semplici molte volte non sono e che l'impazienza del legislatore affastella e complica all'inverosimile -- e propenso piuttosto a ragionare e decidere per principi}''}\footnote{H. Simonetti, \emph{Brevi note sul possibile uso dell'intelligenza artificiale nel processo amministrativo,} cit. p.~931.}\emph{.} Viene allora spontaneo asserire che la società umana insiste, a fasi alterne, nella ricerca di meccanismi automatizzati che possano rendere più prevedibili, o quantomeno più condivisibili, i provvedimenti giudiziari. L'hardcore del lavoro del giudicante consiste proprio nell'interpretazione del diritto, e nel potere di superare gli orientamenti precedenti in ragione dell'evoluzione culturale della propria comunità di riferimento: su questa attività interpretativa ed innovativa e sulla loro delegabilità all'IA occorre interrogarsi. Nell'ambito dell'ausilio al potere giudiziario, l'algoritmo può essere assimilato a uno strumento tecnico-operativo, analogo a una formula di calcolo o a una griglia valutativa: esso supporta il giudice senza incidere sull'esercizio del potere decisorio in senso proprio. Si sposa qui quell'orientamento dottrinario\footnote{Ivi, H. Simonetti, \emph{Brevi note sul possibile uso dell'intelligenza artificiale nel processo amministrativo}, cit. par. 3: secondo l'autore l'IA va considerata ``uno strumento al servizio, in particolare, del principio dello \emph{jura novit curia}, principio fondamentale che vincola il giudice all'osservanza delle norme di diritto e che lo rende però libero di applicare quelle che meglio ritiene adattabili al caso concreto, anche mutando la qualificazione giuridica proposta dalle parti. Principio che postula una conoscenza del diritto che diamo per scontata trascurando spesso, o fingendo che non vi sia, il problema della conoscenza effettiva di un diritto divenuto sempre più esteso e sempre più instabile. Tanto più nel campo del diritto amministrativo dove manca un codice di diritto sostanziale e dove si è al cospetto di una pluralità di fonti, di livelli normativi diversi, di una legislazione speciale in continua evoluzione''.} che vede l'IA come una delle modalità di implementazione dello iura novit curia, ossia un mezzo concreto per il giudice di conoscere seriamente tutto il diritto, specialmente in quelle materie, come il diritto amministrativo, ove ciò è quantitativamente impossibile. L'algoritmo non esprime alcuna forma di autonomia cognitiva, né è dotato di capacità adattiva: la relazione tra \emph{input} e \emph{output} è rigidamente strutturata e, almeno in astratto, integralmente ricostruibile\footnote{W. Barfield, \emph{The Cambridge Handbook of the Law of Algorithms}, \emph{An Introduction to Law and Algorithms}, Cambridge, 2020, pp.~3 - 15.}. L'attività interpretativa, usata o abusata che sia\footnote{A onor del vero la dottrina recente rileva più l'abuso che l'uso: F. Saitta, \emph{Interprete senza spartito? Saggio critico sulla discrezionalità del giudice amministrativo}, Napoli, ESI, 2023.; A. Panzarola, \emph{Principi e regole in epoca di utilitarismo processuale}, Bari, 2022.}, non contraddistingue esclusivamente la giurisdizione amministrativa, e ad onor del vero nemmeno la giurisdizione in generale. L'ermeneutica è propria dei magistrati, dei professionisti legali, dei funzionari amministrativi, e finanche del legislatore stesso. Tra le diverse manifestazioni dell'attività interpretativa assume particolare rilievo quella che conduce al superamento di precedenti orientamenti giurisprudenziali. In tale prospettiva si colloca l'\emph{overruling} (o \emph{revirement}), che consiste nell'abbandono di un precedente arresto interpretativo e nell'adozione di una diversa soluzione ermeneutica. La risposta al quesito sulla sostituzione del giudice in questa attività è negativa almeno nel suo senso assoluto, anche grazie ad una preclusione di rango costituzionale come ricostruita in dottrina,\footnote{H. Simonetti, \emph{Brevi note sul possibile uso dell'intelligenza artificiale nel processo amministrativo}, cit. p.~933: \emph{``\emph{Molte, e tutte molto serie, sono le questioni e i problemi di diritto costituzionale che la prospettiva di una decisione giudiziaria robotica solleva. L'idea di fondo è che la Costituzione, in una trama di articoli che vanno dal 25 al 101, comma 2, e al 111, presupponga la giurisdizione come attività eminentemente umana e non ammetta che un uso limitato e solo servente dell'algoritmo}''.}} e alla presa di posizione del legislatore europeo\footnote{AI Act, Regolamento UE 2024/1689, considerando 61: ``L\textquotesingle utilizzo di strumenti di IA può fornire sostegno al potere decisionale dei giudici o all\textquotesingle indipendenza del potere giudiziario, ma non dovrebbe sostituirlo: il processo decisionale finale deve rimanere un\textquotesingle attività a guida umana''.}. L'IA, come rilevato in precedenza, è in grado di classificare casi, individuare pattern ricorrenti e formulare previsioni sulla base di dati e precedenti disponibili; non può, tuttavia, creare diritto, innovare autonomamente la giurisprudenza o assumere decisioni che implicano valutazioni assiologiche e scelte di politica del diritto. In altri termini, essa non è in grado di mettere in discussione criticamente i presupposti teorici e interpretativi su cui si fonda il sapere giuridico, né di superare autonomamente gli orientamenti consolidati della dottrina o della giurisprudenza. Detto altrimenti, l'intelligenza artificiale opera entro i confini del patrimonio conoscitivo sul quale è stata addestrata: può riconoscere regolarità e riprodurre schemi decisionali esistenti, ma non possiede la capacità di rimettere in discussione le categorie concettuali che ne costituiscono il fondamento. Se le viene insegnato che il cielo è azzurro, essa potrà identificare e applicare tale informazione in una molteplicità di contesti, ma non potrà autonomamente interrogarsi sul significato di ``cielo'' o di ``azzurro'', né elaborare una diversa concezione di tali nozioni. In ragione di quanto evidenziato nel secondo paragrafo di questo lavoro, deve ritenersi preclusa allo strumento informatico quella funzione creativa che, entro determinati limiti, caratterizza il diritto di formazione giurisprudenziale. L'intelligenza artificiale, infatti, può riprodurre e sistematizzare orientamenti esistenti, ma non è in grado di contribuire autonomamente alla loro evoluzione attraverso quel processo interpretativo e innovativo che ha storicamente accompagnato l'attività delle autorità giudiziarie\footnote{Sul ruolo creativo dello \emph{ius praetorium} nell'antica Roma si annoverano, tra i contributi classici, O. Lenel, \emph{Das edictum perpetuum}, 1883, e S. Riccobono, \emph{Corso di diritto romano. Formazione e sviluppo del diritto romano dalle XII tavole a Giustiniano. Parte II -- Dalla lex Aebutia a Giustiniano. Ius novum}, Milano, 1933-1934. Quanto alla manualistica, costituiscono riferimenti imprescindibili M. Talamanca, \emph{Lineamenti di storia del diritto romano}, Milano, 1989, e D. Mantovani, \emph{Le formule del processo privato romano. Per la didattica delle istituzioni di diritto romano}, Padova, 2007. Per quanto riguarda le più recenti riflessioni sulle ``decisioni pretoriae'' dell'organo di giustizia costituzionale, si segnala P. Zicchittu, \emph{Il crepuscolo della discrezionalità. Le decisioni della Corte costituzionale di fronte alle scelte politiche del Parlamento}, in \emph{Archivio giuridico online}, IV, n.~1, 2025, pp.~775-836.}.

\subsection{L'overruling automatizzato}\label{loverruling-automatizzato}

L'overruling presuppone la capacità di discostarsi da un orientamento giurisprudenziale consolidato alla luce di una nuova valutazione interpretativa e assiologica. In tali ipotesi, l'intelligenza artificiale incontra un limite strutturale, poiché il superamento del precedente implica una sua messa in discussione critica che eccede la mera rielaborazione dei dati su cui il sistema è stato addestrato. Ne consegue che l'IA potrebbe operare solo nei casi in cui non sia necessario alcun revirement; tuttavia, non è dato stabilire ex ante se una controversia richieda o meno un overruling, giacché tale valutazione è essa stessa il frutto dell'attività interpretativa del giudice. Si configura così un paradosso circolare: l'intelligenza artificiale potrebbe sostituire il magistrato solo quando non occorra un \emph{overruling}, ma la verifica della sua necessità presuppone comunque l'intervento del magistrato\footnote{A ciò si aggiunge il fatto che, secondo la storica filosofia del diritto del realismo americano, qualsiasi pronuncia giudiziaria è in realtà un momento di creazione del diritto, e non di mera applicazione di esso. Celebre opera sui processi mentali del giudice è J. Frank, \emph{Law and the Modern Mind}, New Brunswick and London, 2009 (ed.~orig. 1930).}. In tali termini, l'intelligenza artificiale non potrebbe sostituire il giudice nemmeno nelle controversie che non richiedono soluzioni innovative, poiché non è comunque predeterminabile ex ante la natura della decisione da adottare nel caso concreto. La qualificazione di una fattispecie come ``consolidata'' o ``aperta'' è, infatti, essa stessa il risultato di un'attività interpretativa. Un ipotetico ``\emph{overruling} automatizzato'' potrebbe configurarsi solo come effetto di un intervento esterno sui criteri decisionali del sistema, attraverso la modifica dei dati o dei parametri di addestramento, così da orientarne l'output verso nuove gerarchie di valori e interessi. In tal senso, più che di \emph{overruling} in senso proprio, si tratterebbe di un aggiornamento del modello, analogo a un costante refining dei dati di input, come avviene nei sistemi di navigazione satellitare che si adeguano alle modifiche della viabilità senza possedere alcuna conoscenza diretta del contesto. Ne deriva che l'\emph{overruling}, in senso tecnico, implica sempre un mutamento ``interno'' e consapevole dell'orientamento precedente da parte dell'organo giudicante. Non vi sarebbe, pertanto, un vero revirement se esso fosse imposto dall'esterno, né se consistesse nella mera applicazione di un diverso dataset proveniente da terzi. Per analogia, un autentico \emph{overruling} presupporrebbe un sistema capace di auto-rivedere liberamente le proprie scelte interpretative, ipotesi che non ricorre nei sistemi di intelligenza artificiale, i quali operano entro schemi determinati dall'addestramento e dai successivi aggiornamenti.

L'IA, salvo introduzioni di dati dell'esterno, non può in alcun modo essere a conoscenza della mutata sensibilità sociale e neppure della nuova produzione della letteratura scientifica su una certa materia\footnote{Queste variabili sono antagonistiche delle esigenze di prevedibilità del sistema capitalistico, come ricordato in M. Luciani, \emph{La decisione giudiziaria robotica}, in \emph{Rivista AIC}, n.~3/2018, spec. pp.~872-893.}. Non è quindi configurabile pertanto alcun \emph{overruling} da parte di un giudice algoritmico, restando questa attività connotata da una riserva di umanità fisiologica e tra l'altro legislativamente sancita\footnote{Art. 15, legge n.~132/2025, qui riportato testualmente per evitare fraintendimenti sulla \emph{voluntas legis}, e lasciare che il lettore apprezzi questa eccezionale previsione normativa: \emph{``\emph{Nei casi di impiego dei sistemi di intelligenza artificiale nell\textquotesingle attività giudiziaria è sempre riservata al magistrato ogni decisione sull\textquotesingle interpretazione e sull\textquotesingle applicazione della legge, sulla valutazione dei fatti e delle prove e sull\textquotesingle adozione dei provvedimenti}''.}}, il che è stato ritenuto meritevole di interesse già nella recente produzione scientifica\footnote{Per una valorizzazione di tale \emph{unicum}, relativo alla c.d. riserva di umanità legislativamente tutelata, si segnala C. De Robbio, \emph{L'intelligenza artificiale e il lavoro del giudice. I rischi per l'autonomia e l'indipendenza della magistratura}, in \emph{Giustizia Insieme}, 2026, par. 5.}. L'IA non può mutare convinzioni, giacché non ha convinzione alcuna; non può cambiare idea, perché non ha proprio nessuna idea\footnote{Uno dei rischi messi in luce nel contributo sopra citato di De Robbio è rappresentato dalla verticalizzazione giurisprudenziale, quale effetto dei criteri gerarchici di classificazione delle decisioni che privilegiano le pronunce dei giudici di legittimità rispetto a quelle di merito. In tale prospettiva, l'IA tenderebbe a recepire in modo prevalente gli orientamenti del Consiglio di Stato e della Corte di cassazione, trascurando aprioristicamente quelli delle corti di merito, in quanto non apicali.Ne deriva che all'IA non può propriamente attribuirsi una autonoma capacità ideativa: le eventuali ``posizioni'' che essa restituisce si risolvono, piuttosto, nella riproduzione degli indirizzi espressi dalle giurisdizioni superiori.}. E se qualche convinzione la avesse, sarebbe nient'altro che la sintesi degli orientamenti maggioritari nel passato, più presenti nei \emph{training data}.\footnote{L'evoluzione del pensiero giuridico sarebbe così compromessa, o finanche cristallizzata, dato che le decisioni dell'IA giudicante sarebbero ancorate solo ad una selezione di orientamenti ``privilegiati'': in A.Garapon, J. Lassegue,\emph{La giustizia digitale}, Bologna, 2021, pp.~184-185, ben si dice che ne fanno le spese ``gli avvocati che perorano soluzioni alternative, o che vogliono far cambiare una giurisprudenza'', in quanto partono da una posizione svantaggiata di minoranza.} Essa offre soluzioni probabili considerando le circostanze concrete, e nulla di più. Non considera né percepisce alcun mutamento sociale, nessuna evoluzione filosofica del pensiero e della cultura umana, né registra alcuna alterazione dei principi di convivenza civile in una comunità. Qualsiasi modifica del codice da parte del programmatore o aggiornamento dei dati nella fase di training si tradurrebbe, in realtà, in un mero aggiornamento eteronomo del sistema, non assimilabile a un autentico \emph{overruling}. Si tratterebbe, infatti, di un cambiamento determinato dall'esterno e non di un mutamento spontaneo della linea decisionale dell'algoritmo, quale potrebbe configurarsi solo in presenza di una reale autonomia interpretativa. In tal senso, la distinzione tra aggiornamento e \emph{overruling} riflette il fatto che quest'ultimo implica una capacità di revisione consapevole delle proprie scelte, propria esclusivamente dell'agire umano, in quanto caratterizzato dall'integrazione tra elaborazione cognitiva e coscienza valutativa\footnote{L'imperitura scissione esistente tra intelligenza artificiale e coscienza è causa primigenia della centralità dell'uomo, e l'IA è prova di questa scissione. Si segnala il celebre Y. N. Harari,\emph{Homo Deus. Breve storia del futuro}, Milano, 2017, p.~484, secondo cui grazie al progresso tecnologico ``l'intelligenza si sta affrancando dalla coscienza''.}.

\subsection{L'IA nel processo, l'IA sotto processo}\label{lia-nel-processo-lia-sotto-processo}

Le considerazioni che precedono evidenziano la crescente complessità delle questioni trattate, nonché la particolare delicatezza che le stesse assumono in sede processuale. L'esercizio della funzione giurisdizionale è già di per sé attraversato da fisiologiche tensioni interpretative e da inevitabili margini di incertezza; in tale contesto l'innovazione tecnologica funge da moltiplicatore potenziale delle medesime problematiche\footnote{Sulle applicazioni dell'IA alla realtà giuridica, v. G. Sartor, \emph{L'intelligenza artificiale e il diritto}, cit., pp.~103 ss.}. Il processo amministrativo emerge qui ``come lo spazio nel quale è possibile giudicare l'intelligenza artificiale e giudicare con l'intelligenza artificiale''\footnote{A. Pajno, \emph{La transizione digitale e il contributo dei giuristi. L'importanza del processo amministrativo}, in M.Ramajoli (a cura di), \emph{Una giustizia amministrativa digitale?}, 2023, p.~15.}\emph{,} fermo restando che l'oggetto dell'indagine non è la valutazione dell'IA in sé, bensì la verifica delle condizioni che ne consentano un impiego compatibile e fruttuoso nell'ambito della disciplina processuale. La giustizia è stata a lungo oracolare\footnote{Il riferimento alla dimensione ``oracolare'' della giustizia trova un significativo antecedente in P. Dawson, \emph{The Oracles of the Law}, Ann Arbor, University of Michigan Press, 1968.}, e non pare conveniente augurarsi l'avvento di un'IA giudicante neo-oracolare\footnote{Sulla giustizia neo-oracolare a mezzo algoritmico, vedi G. Montedoro, \emph{Intelligenza artificiale e digitalizzazione fra Costituzione. amministrazione e giustizia,} in\emph{\href{http://giustizia-amministrativa.it}{\ul{giustizia-amministrativa.it}},} 2024.} per colmare l'incertezza fisiologica dei processi decisionali umani\footnote{Queste caratteristiche intrinseche sono ineliminabili, specialmente nei settori ove resiste \emph{le principe de primauté humaine,}come asserisce il \emph{Conseil d'Etat} nel documento \emph{Intelligence artificielle ed action publique: construir la confiance, servir la performance},2022.}. Tra l'altro, l'uso di questi strumenti nell'attività giudiziaria è già classificato ad alto rischio dalla normative europea\footnote{Sull'AI Act del 2024, v. B. Marchetti, \emph{Intelligenza artificiale, poteri pubblici, rule of law}, in \emph{Rivista italiana di diritto pubblico e comunitario}, 2024, p.~49 ss. Sulle proposte de iure condendo elaborate a partire dal Regolamento europeo, v. G. Gallone, \emph{Riserva di umanità, intelligenza artificiale e funzione giurisdizionale alla luce dell'IA Act. Considerazioni (e qualche proposta) attorno al processo amministrativo che verrà}, in \emph{www.giustizia-amministrativa.it}.}. A prescindere da ogni linea definitiva, sembrano da condividere le parole sottili di un autore sagace, che invertono la prospettiva di osservazione sull'IA: \emph{``\emph{questa non è solo un potenziale oggetto di regolazione da parte del diritto, ma è un sicuro produttore di effetti sul diritto}''}\footnote{M. Luciani, \emph{Può il diritto disciplinare l'Intelligenza Artificiale? Una conversazione preliminare}, in \emph{Bilancio Comunità Persona}, n.~2/2023, p.~15.}. Ed essendo il diritto un fatto umano, è innegabile che sia l'umano a dover esercitare il dominio su esso, senza esternalizzare ad oggetti inorganici le sue funzioni: ciò avviene ammettendo l'IA nel suo uso servente, ma riservando la primazia processuale all'antropomorfo giudicante incaricato di governare la complessità dello Stato costituzionale\footnote{G. Zagrebelsky, \emph{Il diritto mite}, Torino, 1992, p.~213, ove per complessità l'autore intende la ``necessaria, mite coesistenza di legge, diritti e giustizia''.}. Giacché, in caso contrario, sarebbe incommensurabile \emph{``\emph{il rischio che a forza di delegare, e di non fare, non si sappia poi neppure più come controllare il fare, se giusto o sbagliato, degli altri}''}\footnote{H. Simonetti, \emph{Brevi note sul possibile uso dell\textquotesingle intelligenza artificiale nel processo amministrativo}, cit., p.~934.}.

\section{Processo amministrativo e intelligenza artificiale. Discrezionalità giudiziale e discrezionalità amministrativa a confronto}\label{processo-amministrativo-e-intelligenza-artificiale.-discrezionalituxe0-giudiziale-e-discrezionalituxe0-amministrativa-a-confronto}

Dopo aver esaminato il rapporto tra funzione giurisdizionale e intelligenza artificiale, con particolare riferimento all'attività ermeneutica, occorre ora soffermarsi su un ulteriore profilo, forse ancor più problematico: l'interazione tra intelligenza artificiale e discrezionalità giudiziale - intesa come attività di valutazione e ponderazione tra interessi contrapposti - che connota in modo peculiare il processo amministrativo. È in questo ambito che emergono le principali criticità, tanto sul piano teorico quanto su quello applicativo, poiché l'esercizio della discrezionalità implica, non di rado, scelte fondate sulla comparazione tra principi fondamentali dell'ordinamento e, talvolta, tra interessi non previamente determinati dal legislatore, come accade in modo paradigmatico nella discrezionalità cautelare. Le decisioni cautelari, infatti, non si esauriscono in una valutazione centrata sulla sola posizione soggettiva del ricorrente, ma richiedono un bilanciamento che tenga conto anche dell'impatto della misura sull'effettivo perseguimento delle finalità pubbliche\footnote{La dimensione cautelare del processo amministrativo si caratterizza per un'elevata discrezionalità giudiziale, che trova il suo fondamento nella necessità di operare un bilanciamento tra interessi eterogenei, non rigidamente predeterminato dal legislatore. In tale prospettiva, la dottrina ha da tempo evidenziato come il giudizio cautelare non possa essere ridotto a una verifica meramente tecnica dei presupposti del \emph{fumus boni iuris} e del \emph{periculum in mora}, ma implichi una più ampia valutazione comparativa degli interessi coinvolti.In questa direzione, A. Romano Tassone (\emph{Risarcibilità del danno e tutela cautelare amministrativa}, in \emph{Dir. amm.}, 2001, p.~26 ss.) ha proposto una nota classificazione delle diverse concezioni dell'interesse pubblico rilevante nella fase cautelare, distinguendo tra: a) l'interesse generale alla legalità dell'azione amministrativa; b) l'interesse concreto perseguito dall'amministrazione nel caso specifico; c) l'interesse generale alla corretta composizione degli interessi pubblici e privati coinvolti. Tale tripartizione conduce l'Autore a sottolineare il carattere prevalentemente oggettivo della tutela amministrativa, e in particolare di quella cautelare, che trascende la mera protezione della posizione soggettiva del ricorrente. Più recentemente, A. Travi (\emph{Lezioni di giustizia amministrativa}, Torino, 2021, p.~294) evidenzia come il giudice cautelare sia chiamato a svolgere una valutazione comparata di tutti gli interessi in gioco, secondo criteri non puntualmente codificati, rimessi al suo prudente apprezzamento. Analogamente, G. Leone (\emph{Elementi di diritto processuale amministrativo}, Padova, 2019, p.~240) sottolinea che il giudice deve considerare non solo il pregiudizio del ricorrente, ma anche quello che la misura cautelare potrebbe arrecare all'amministrazione e alla collettività, nonché procedere, ove necessario, a un confronto tra interessi pubblici tra loro confliggenti.Infine, M. Lipari (\emph{La nuova tutela cautelare degli interessi legittimi: il ``rito appalti'' e le esigenze imperative di interesse generale}, in \emph{www.federalismi.it}, 5/2017, p.~17) rimarca come la decisione cautelare non possa essere confinata alla prospettiva unilaterale del ricorrente, ma debba necessariamente tener conto, in modo unitario, della complessa rete di interessi giuridici coinvolti, attribuendo alle finalità pubbliche perseguite dall'amministrazione un ruolo essenziale nel giudizio.Nel loro complesso, tali contributi convergono nel delineare una concezione della discrezionalità giudiziale cautelare come attività di bilanciamento strutturalmente aperta e relazionale, nella quale la tutela dell'interesse individuale si intreccia inscindibilmente con la salvaguardia degli interessi pubblici.}. In tali ipotesi, l'attività del giudice trascende l'interpretazione in senso stretto, traducendosi in un giudizio composito, nel quale si intrecciano valutazioni normative, accertamenti fattuali e apprezzamenti di carattere assiologico. Alla luce di tali premesse, appare problematico ipotizzare l'impiego di strumenti di intelligenza artificiale nel momento ponderativo della fase cautelare del processo amministrativo. In questo segmento, infatti, la discrezionalità giudiziale si manifesta con particolare intensità, imponendo al giudice di operare in contesti connotati da elevata incertezza e da una marcata variabilità delle circostanze rilevanti. Essa implica valutazioni non rigidamente predeterminate, rimanendo strutturalmente orientata dal principio di effettività della tutela giurisdizionale, senza che ciò ne implichi una necessaria e costante prevalenza. È frequente, ad esempio, che le procedure espropriative non vengano sospese in sede cautelare, in considerazione della necessità di assicurare la tempestiva realizzazione di opere pubbliche strategiche, destinate a soddisfare, in tempi brevi - talvolta strettissimi - esigenze di particolare rilievo per la collettività. Inoltre, soprattutto l'esperienza maturata durante l'emergenza da Covid-19 ha evidenziato come il giudice sia chiamato a svolgere valutazioni complesse e di natura qualitativa, finalizzate al bilanciamento tra interessi pubblici e privati tra loro confliggenti, non sempre predeterminati dal legislatore, ma desumibili da principi generali e valori dell'ordinamento.\\
In tale contesto, la sede cautelare assume un ruolo centrale, poiché in essa la controversia viene spesso definita in modo di fatto irreversibile, soprattutto nei casi in cui si impugnano provvedimenti a efficacia temporanea, come quelli adottati durante l'emergenza pandemica (ad esempio, le misure di sospensione delle attività scolastiche). Il potere cautelare si configura quale strumento attraverso cui vengono apprezzate le possibili ricadute della misura sia sull'azione amministrativa sia sugli interessi pubblici sottesi al provvedimento impugnato. Si tratta di un potere particolarmente ampio, esercitato in assenza di criteri legali di comparazione rigidamente predeterminati e di una gerarchia definita tra gli interessi coinvolti. Ne consegue che il giudice è chiamato a individuare, caso per caso, il punto di equilibrio più adeguato, mediante un giudizio prudenziale che integra in modo unitario elementi normativi, fattuali e assiologici e che, proprio per tale complessità, risulta difficilmente riducibile a procedure meccaniche o a modelli di calcolo probabilistico. Una diversa impostazione\footnote{Sostenuta, tra gli altri, da G. D'Elia, \emph{Giudice umano e Intelligenza artificiale: il processo amministrativo alla prova dell'imparzialità nell'era delle decisioni automatizzate}, in \emph{PA, Persona e amministrazione}, 2024, p.~1267 ss., spec. p.~1281.} ipotizza che sistemi di intelligenza artificiale, addestrati mediante algoritmi o tecniche di \emph{deep learning}, possano riprodurre decisioni complesse, simulando il giudizio umano attraverso l'elaborazione di input adeguatamente strutturati. Tuttavia, tale prospettiva non considera che la discrezionalità giudiziale presenta un nucleo intrinsecamente non formalizzabile, in quanto implica la comparazione tra principi generali - anche di rango costituzionale - spesso privi di una gerarchia definita. Nelle decisioni cautelari, infatti, il giudice non si limita a classificare dati secondo logiche probabilistiche, ma interpreta norme, integra informazioni incomplete o incerte\footnote{Nel processo amministrativo vige il principio dispositivo attenuato dal metodo acquisitivo, che consente al giudice amministrativo di acquisire d'ufficio materiale probatorio in relazione ai fatti per i quali sia stato fornito un principio di prova.}, valuta le conseguenze pratiche delle proprie decisioni e pondera le priorità tra interessi concorrenti, operando scelte prudenziali che presuppongono sensibilità etica e capacità di giudizio contestuale, nonché l'attenzione ai mutamenti della sensibilità sociale e del contesto valoriale di riferimento, elementi che difficilmente possono essere intercettati e rielaborati da sistemi di intelligenza artificiale. Attribuire all'intelligenza artificiale il potere di svolgere simili valutazioni significherebbe, pertanto, sostituire valutazioni soggettive proprie dell'esperienza umana con valutazioni parimenti non neutrali, ma elaborate da un sistema automatizzato\footnote{Sulla non neutralità dell'intelligenza artificiale v., ex multis, Cathy O\textquotesingle Neil, \emph{Weapons of Math Destruction. How Big Data IncreasesInequality and Threatens Democracy}, New York, 2016.}. Un analogo ordine di considerazioni si estende al giudizio speciale in materia di contratti pubblici, nel quale il giudice amministrativo è chiamato a bilanciare l'interesse alla tutela delle situazioni giuridiche soggettive - spesso coincidente con la salvaguardia della concorrenza\footnote{Sul rapporto di strumentalità esistente tra salvaguardia dell'interesse dell'operatore economico a conseguire il contratto e tutela della concorrenza v. E. Sticchi Damiani, \emph{Annullamento dell'aggiudicazione e inefficacia funzionale del contratto}, in \emph{Dir. proc. amm}., 2011, p.~262, secondo il quale non può dubitarsi che «la tutela della concorrenza si determini, sotto il profilo dei contenuti, nella misura più piena non solo allorché un'aggiudicazione illegittima faccia venir meno il contratto, ma soprattutto allorché in quell'aggiudicazione subentri chi aveva titolo a conseguirla così che la pronuncia giurisdizionale produca il massimo dei suoi effetti ricostruttivi».} - con l'esigenza di assicurare la continuità nell'erogazione dei servizi e la tempestiva realizzazione delle opere pubbliche. La tensione tra tali esigenze emerge in modo particolarmente evidente nella disciplina di cui all'art. 122 c.p.a., ove il legislatore delinea una sostanziale equiordinazione tra l'effettività della tutela giurisdizionale, da un lato, e l'interesse pubblico alla celere esecuzione delle prestazioni contrattuali, dall'altro. In questo caso, tuttavia, gli interessi rilevanti ai fini della decisione risultano espressamente individuati a livello legislativo, con la conseguenza che l'attività di bilanciamento del giudice si svolge entro un perimetro normativamente tracciato, risultando più vicina, per struttura, a quella esercitata in sede amministrativa. In questo contesto, il giudice può dichiarare l'inefficacia del contratto concluso all'esito di una procedura illegittima, valorizzando le esigenze di tutela giurisdizionale e di concorrenza, oppure, in alternativa, ritenere prevalente l'interesse pubblico alla continuità del servizio o al completamento dell'opera, mantenendo in vita il rapporto contrattuale nonostante i vizi che ne abbiano inficiato la formazione. Ne consegue che, in tali ipotesi, i poteri del giudice travalicano i confini dell'interpretazione in senso stretto, estendendosi a valutazioni di opportunità in concreto che implicano ampi margini di apprezzamento. In questa prospettiva, è stato efficacemente osservato come, nella ponderazione degli interessi contrapposti - segnatamente con riguardo alle sorti del contratto - il giudice tenda a esercitare una discrezionalità che si avvicina, per intensità e struttura, a quella propria dell'azione amministrativa\footnote{F. Saitta, \emph{Interprete senza spartito? Saggio critico sulla discrezionalità del giudice amministrativo}, Napoli, 2023, p.~432 ss., che ha espresso perplessità in ordine alla scelta legislativa di affidare al giudice amministrativo un bilanciamento così delicato, suscettibile di incidere sull'equilibrio tra funzione giurisdizionale e funzione amministrativa}. Tuttavia, a differenza di quanto avviene nelle decisioni cautelari, tale discrezionalità si esercita su interessi previamente selezionati dal legislatore, e non su un novero aperto e indeterminato di valori. È in queste situazioni che si manifesta la complessità della discrezionalità giudiziale e il carattere insostituibile del giudizio umano, che non può essere replicato integralmente da algoritmi o procedure automatizzate. Anche gli strumenti di intelligenza artificiale più avanzati risultano infatti incapaci di cogliere appieno la specificità di un singolo caso, poiché mancano della visione olistica e contestualizzata propria dell'essere umano. In altre parole, se i dati utilizzati per l'addestramento del modello non riflettono la varietà dei casi reali o non includono informazioni relative a circostanze individuali o situazioni particolari, l'algoritmo rischia di produrre valutazioni parziali o inadeguate rispetto alla singola situazione concreta, compromettendo così l'effettività della tutela giurisdizionale. Il confronto tra discrezionalità giudiziale e discrezionalità amministrativa mette in luce una differenza strutturale di fondo. La discrezionalità amministrativa si configura come attività funzionalizzata: il legislatore individua gli interessi rilevanti e orienta l'azione amministrativa verso il perseguimento di finalità predeterminate. Ad esempio, nell'esercizio del potere di autotutela, l'amministrazione è chiamata a bilanciare, da un lato, l'interesse pubblico al ripristino della legalità violata e, dall'altro, l'interesse del privato alla conservazione del vantaggio conseguito. In tale valutazione comparativa, essa deve attenersi anche ai criteri indicati dal legislatore, quali il decorso del tempo e la tutela dell'affidamento del destinatario del provvedimento. Ne deriva che la discrezionalità amministrativa si esercita nella scelta tra opzioni già delineate dalla norma, presentando un grado relativamente elevato di predeterminazione che, almeno in linea teorica, rende più agevole il ricorso a strumenti di intelligenza artificiale in funzione di supporto decisionale\textsuperscript{.} Anche la discrezionalità giudiziale implica un potere di scelta tra diverse soluzioni possibili: si pensi, ad esempio, alla decisione di accogliere o rigettare una domanda cautelare, sulla base della valutazione comparativa tra gli interessi coinvolti. Tuttavia, nella discrezionalità giudiziale - a differenza di quella amministrativa - può risultare assente quello che può definirsi il ``momento prefigurativo'', ossia la previa tipizzazione legislativa degli interessi da perseguire. Tale profilo emerge con particolare evidenza nelle decisioni cautelari, nelle quali il giudice è chiamato a operare un bilanciamento tra interessi non previamente predeterminati dal legislatore, ma desumibili dai principi costituzionali e dai valori generali dell'ordinamento. L'art. 55 c.p.a., infatti, non individua puntualmente i criteri del bilanciamento né impone la considerazione espressa dell'interesse pubblico, a differenza di quanto previsto dall'art. 120, comma 10, c.p.a. Ne consegue che la discrezionalità giudiziale esercitata in sede cautelare non è sempre riconducibile a un modello funzionalizzato: il giudice è chiamato a svolgere un'attività di ponderazione nella quale gli interessi rilevanti non sono tipizzati e il parametro decisionale deve essere costruito caso per caso, anche mediante il ricorso a principi generali e valori di ampia portata. In questa prospettiva, la discrezionalità giudiziale non si esaurisce in una mera comparazione tra interessi normativamente predeterminati, ma implica anche l'individuazione degli interessi rilevanti e dei criteri del loro bilanciamento, non previamente definiti dal legislatore. Ne deriva che l'impiego dell'intelligenza artificiale incontra limiti strutturali più marcati nell'ambito della discrezionalità giudiziale. Le decisioni del giudice amministrativo implicano, infatti, un'operazione unitaria che combina interpretazione normativa contestuale, integrazione di informazioni incomplete o incerte e ponderazione tra principi e interessi spesso privi di una gerarchia definita. Tali decisioni sono strutturalmente orientate al principio di effettività della tutela giurisdizionale, quale fine primario del processo amministrativo, ma richiedono altresì la considerazione degli interessi generali della collettività e dell'interesse pubblico sotteso al provvedimento impugnato e, dunque, del risultato amministrativo. Il giudice, pertanto, non può risolversi in scelte meramente soggettive o strumentali, ma è tenuto a garantire una tutela concreta delle situazioni giuridiche, attraverso un bilanciamento fondato su criteri di ragionevolezza, proporzionalità e prudenza. Ne consegue che l'esercizio della discrezionalità giudiziale - specie quando non sia ancorato a interessi legislativamente predeterminati - appare inscindibilmente legato all'apporto umano\footnote{È pur vero che, secondo la cosiddetta legge dei ritorni acceleranti descritta da Ray Kurzweil (The Age of Spiritual Machines, 1999), la tecnologia computazionale è destinata a progredire a un ritmo esponenziale, fino a superare le capacità della specie umana. In tale scenario, le macchine saranno in grado di elaborare quantità di dati inaccessibili al cervello umano, con costi energetici inferiori e maggiore efficienza; la loro capacità di calcolo potrebbe divenire così elevata da rendere difficile comprendere e valutare pienamente i risultati da esse prodotti. Una simile eventualità implicherebbe, tuttavia, il rischio di un giudizio ``quasi oracolare'', nel quale l'uomo perderebbe il controllo interpretativo sui processi decisionali automatizzati. Lo stesso Kurzweil prospetta, inoltre, che, intorno al 2045, gli esseri umani possano integrarsi con le macchine, potenziando le proprie capacità cognitive mediante interfacce cervello-computer, mentre le macchine non biologiche diverrebbero capaci di auto-migliorarsi in modo continuo, fino a rendere il loro funzionamento difficilmente comprensibile per gli esseri umani non potenziati. Tutto ciò riconduce alla domanda iniziale (v., \emph{supra}, nt. 5): è effettivamente auspicabile sostituire il giudice umano con sistemi di intelligenza artificiale, anche laddove questi ultimi fossero in grado di garantire risultati apparentemente migliori in termini di giustizia sostanziale, equità delle decisioni e adeguatezza rispetto alla complessità degli interessi coinvolti?}. La prospettata sostituzione del giudice amministrativo con sistemi automatizzati, talora evocata come soluzione a inefficienze del sistema, rischierebbe di generare criticità ben più gravi di quelle che intende risolvere. L'intelligenza artificiale non può essere interposta come un filtro opaco tra il giudice e tale complessa attività: qualsiasi strumento che pretenda di valutare, interpretare o ponderare in sua vece finirebbe per espropriare il magistrato del nucleo essenziale della sua funzione.

\section{Conclusioni: discrezionalità giudiziale e limiti dell'automazione decisionale}\label{conclusioni-discrezionalituxe0-giudiziale-e-limiti-dellautomazione-decisionale}

Alla luce delle considerazioni svolte, il rapporto tra intelligenza artificiale e funzione giurisdizionale non appare riconducibile a una logica sostitutiva del decisore umano, dovendo piuttosto essere inquadrato entro una dimensione di ausilio tecnico-operativo\footnote{Sulla capacità dell'algoritmo di uniformare giustizia, equità e legalità vedi B. Romano, \emph{Algoritmi al potere. Calcolo giudizio pensiero}, Torino, 2018, p.~27 ss. Secondo l'Autore, questi tre concetti sono tradizionalmente autonomi e separati, e solo il Giudice poteva unirle casisticamente in modo armonico. L'algoritmo può, come prima faceva solo il magistrato umano, condurre ad unità queste tre dimensioni del processo.}, idonea a incrementare l'efficienza, favorire l'uniformità applicativa e rafforzare la capacità predittiva del sistema, senza incidere sul nucleo propriamente valutativo della decisione, che resta affidato alla responsabilità del giudice. Le principali criticità emergono, anzitutto, con riguardo all'esercizio della discrezionalità amministrativa, ambito nel quale l'intervento algoritmico incontra limiti strutturali significativi. In tale quadro assume rilievo anche il principio di legalità, nella misura in cui è il legislatore ad attribuire alla pubblica amministrazione - e non a sistemi di intelligenza artificiale - la titolarità e l'esercizio del potere, soprattutto laddove si tratti di potestà autoritative implicanti scelte discrezionali e valutazioni comparativo-ponderative\footnote{Sull'ascesa e i limiti della legalità algoritmica si veda E. Carloni, \emph{I principi della legalità algoritmica. Le decisioni automatizzate di fronte al giudice amministrativo}, cit., pp.~271 ss.; Id., \emph{Critica dell'amministrazione artificiale}, cit., spec. pp.~91 ss. In argomento, v. anche L. Torchia, \emph{Lo stato digitale}, cit. p.~129 ss.}. Ne consegue che l'eventuale apporto algoritmico non può che collocarsi in una dimensione ancillare e strumentale, restando preclusa ogni forma di traslazione della decisione in capo al sistema automatizzato. Ciò nondimeno, non può escludersi, in via generale, un impiego dell'intelligenza artificiale in funzione propositiva del provvedimento, quale strumento capace di elaborare ipotesi decisionali suscettibili di essere vagliate e, se del caso, rielaborate dall'autorità procedente. In questa prospettiva, autorevole dottrina valorizza la potenziale attitudine dei sistemi di intelligenza artificiale a concorrere, almeno in via astratta, all'esercizio di attività discrezionali. Si è, infatti, efficacemente osservato come le tecniche di intelligenza artificiale, e segnatamente le reti neurali, talora qualificate come ``scienza del giusto peso'', possano realizzare sofisticate operazioni di bilanciamento tra interessi e valori, pervenendo, in taluni casi, a esiti caratterizzati da un apprezzabile grado di coerenza interna\footnote{F. Romeo, \emph{Algoritmi di giustizia ed equità nel diritto}, in \emph{i-lex}, 2021, p.~12.}. Tale considerazione acquista particolare rilievo nelle ipotesi in cui gli interessi pubblici contrapposti risultino previamente tipizzati dal legislatore, come accade nell'ambito delle norme attributive di poteri discrezionali alla pubblica amministrazione. In tali contesti, infatti, l'attività di ponderazione si svolge entro una cornice almeno parzialmente definita, rispetto alla quale l'apporto algoritmico può rivelarsi più incisivo, pur senza esaurire la complessità della decisione. Trasposta sul piano processuale, l'intelligenza artificiale mostra già oggi una significativa utilità nello svolgimento di compiti materiali e organizzativi, quali la gestione dei ruoli d'udienza, la calendarizzazione dei procedimenti e la ricerca avanzata dei precedenti giurisprudenziali, nonché, con maggiori criticità\footnote{Le criticità derivano dalla natura probabilistico-induttiva dei sistemi di \emph{machine learning}, che, nel ricondurre la fattispecie a precedenti consolidati, possono determinare una sovra-estensione dei pattern decisionali già esistenti, qualificando come ``non complessa'' una controversia che presenti invece elementi di discontinuità rispetto ai casi precedentemente decisi. Tali criticità, tuttavia, potrebbero risultare in parte attenuate dal progressivo sviluppo dei modelli di addestramento e dall'incremento delle capacità computazionali, che consentirebbero una più raffinata modellazione delle categorie decisionali, pur senza eliminarne integralmente la portata.}, nell'attività di disvelamento delle cause prive di particolare complessità ai sensi dell'art. 72-\emph{bis} c.p.a., contribuendo così a migliorare l'efficienza complessiva del sistema senza incidere direttamente sul contenuto della decisione. Permangono, tuttavia, rilevanti perplessità con riferimento all'applicazione dell'intelligenza artificiale all'attività interpretativa del giudice, specie ove si consideri il fenomeno dei \emph{revirement} giurisprudenziali, espressione della fisiologica evoluzione del diritto vivente\footnote{Sui rischi di un neopositivismo algoritmico, v. L. Torchia, \emph{Lo Stato digitale}, cit., p.~199, la quale segnala, in particolare, il pericolo di rafforzare il conformismo decisionale attraverso l'affidamento a sistemi automatizzati che tendono a standardizzare le soluzioni giuridiche. In tale prospettiva, si evidenzia altresì il rischio di ``ammantare di una (apparente) oggettività'' le decisioni del giudice, le quali finiscono per essere percepite come il prodotto neutrale di procedure algoritmiche, piuttosto che come esito di un'attività interpretativa e valutativa connotata da discrezionalità e responsabilità.}. Tale dinamica, lungi dall'essere riducibile a modelli predittivi stabilizzati, si radica in una dimensione propriamente antropica, nella quale assumono rilievo la sensibilità istituzionale del giudice e la sua capacità di rielaborazione critica del dato normativo, anche alla luce dei mutamenti della sensibilità sociale e del contesto ordinamentale di riferimento, che ne orientano l'interpretazione in senso evolutivo\footnote{Questo approccio rende più appropriato parlare, con riferimento all'impiego di tali tecnologie nella giustizia amministrativa italiana, di ``intelligenza accelerata'' piuttosto che di ``intelligenza artificiale'', così valorizzandone la funzione meramente ausiliaria rispetto all'attività decisoria umana: B. Bruno e D.F. Sivilli, \emph{Digitalizzazione della Giustizia amministrativa italiana: l'impiego delle tecnologie di IA}, in \emph{Il Processo}, n.~3/2024, p.~775 ss.}. Simili criticità, già ravvisabili sul piano dell'attività interpretativa, emergono tuttavia con maggiore evidenza con riguardo all'esercizio della discrezionalità giudiziale. Essa si configura, infatti, quale attività di bilanciamento tra principi e interessi non rigidamente gerarchizzati, esercitata caso per caso attraverso un giudizio prudenziale che integra elementi normativi, fattuali e assiologici. In sede cautelare, poi, tale carattere si accentua ulteriormente, poiché gli interessi rilevanti non risultano integralmente predeterminati, ma vengono individuati dal giudice alla luce dei principi generali dell'ordinamento. In questa prospettiva, come efficacemente rilevato, ciò ``lascia al giudice ampia discrezionalità, aprendo grandi brecce all'ingresso di valori innominati'' \footnote{F. Romeo, \emph{Algoritmi di giustizia ed equità nel diritto}, cit., p.~13.}, non sempre riconducibili all'ordinamento, ma talvolta espressione della sensibilità del singolo giudicante. Il rilievo evidenzia, da un lato, la fisiologica apertura del giudizio ai valori e, dall'altro, la necessità di un controllo argomentativo particolarmente rigoroso, che solo una motivazione piena e trasparente è in grado di assicurare. Nondimeno, proprio il carattere assiologico e contestuale del giudizio giurisdizionale - che implica selezione, ponderazione e giustificazione dei valori in gioco - impedisce di attribuire all'intelligenza artificiale un ruolo sostitutivo. L'algoritmo può, al più, concorrere alla ricognizione degli interessi rilevanti e alla simulazione di possibili esiti decisionali; ma il momento della scelta, e soprattutto quello della sua giustificazione, restano inscindibilmente affidati al giudice. Ne discende che una decisione orientata ai valori esige un onere motivazionale particolarmente intenso: solo attraverso una motivazione completa, intellegibile e contro-argomentabile è possibile evitare che la discrezionalità si traduca in opacità o arbitrio, preservando la funzione di garanzia propria della giurisdizione. In conclusione, la prospettiva di una giurisdizione integralmente automatizzata appare, allo stato, non solo problematica sul piano operativo, ma anche concettualmente incompatibile con la struttura stessa della discrezionalità giudiziale, la quale resta intrinsecamente connessa all'apporto umano, quale luogo di sintesi tra razionalità ed esperienza, tra regola e valore, nonché presidio ultimo di legittimazione della decisione giuridica. Il problema, peraltro, non attiene soltanto alla possibilità tecnica della decisione automatizzata, quanto piuttosto alla sua accettabilità rispetto al valore della giustizia. L'idea di una produzione seriale di decisioni ad opera di un'\,``IA giudicante'' evoca scenari distopici, più che rappresentare un obiettivo auspicabile\footnote{Sui rimedi a una possibile deriva produttivistica, si segnala ancora C. De Robbio, \emph{L'intelligenza artificiale e il lavoro del giudice. I rischi per l'autonomia e l'indipendenza della magistratura}, cit., par. 5.}. Un atteggiamento di acritica tecnofilia\footnote{Per un'analisi dell'endiade tecnofobia e tecnofilia, con alcuni cenni sull'origine dei termini, vedi R. Barry, \emph{Technophobia and Technophilia}, in \emph{British Journal of Psychotherapy}, 1993, spec. pp.~188-195; T. Edward, \emph{Confessions Of A Technophile}, Raritan, 2002, spec. p 135.}rischia, infatti, di tradursi in una forma di ottimismo ingenuo, incapace di cogliere la complessità del fenomeno giuridico. In tale prospettiva, appare ancora attuale l'insegnamento di chi ha messo in guardia rispetto ai limiti dei progetti razionalizzatori di matrice altomodernista\footnote{Sul macrotema del modernismo e sui suoi tratti identitari si vedano P.J. Taylor, \emph{Modernities: A GeohistoricalInterpretation}, Minneapolis, 1999; T. Ruprecht, \emph{Socialist High Modernity and Global Stagnation: A Shared History of Brazil and the Soviet Union During the Cold War}, in \emph{Journal of Global History}, 6, n.~3/2011, p.~505 ss.}, evidenziando come l'aspirazione a uniformare e semplificare la realtà sociale finisca spesso per trascurarne le irriducibili specificità. Analogamente, l'idea - radicata nella tradizione giuridica post-rivoluzionaria\footnote{J.C. Scott, \emph{Authoritarian High Modernism}, in S. Fainstein, J. DeFilippis (a cura di), \emph{Readings in Planning Theory}, Chichester, 2015, p.~75 ss.} - di una piena omogeneità dei consociati e delle fattispecie normative si riflette nelle diverse branche del diritto, ma si scontra, nella prassi applicativa, con la varietà e la complessità delle situazioni concrete, che esigono un'opera interpretativa e valutativa non integralmente delegabile alla macchina. È proprio qui che si manifesta uno dei più frequenti limiti del legislatore: l'ipersemplificazione di concetti, entità e soggetti conduce a una rappresentazione riduttiva della realtà sensibile. Si finisce così per osservare il Sole e la Luna attraverso i medesimi occhiali scuri, concepiti per attenuare l'abbagliante luce del primo ma utilizzati, impropriamente, anche dinanzi alla tenue luminosità della seconda. Ne deriva una percezione alterata delle differenze che caratterizzano il reale e che, invece, il diritto è chiamato a riconoscere e valorizzare. Non può sostenersi, in termini assoluti, che la discrezionalità giudiziale rientri tra le attività integralmente replicabili da un sistema di intelligenza artificiale. Essa, infatti, per sua stessa natura, si sottrae a una completa formalizzazione: non è interamente traducibile in regole predeterminate, né può essere pienamente standardizzata o riprodotta mediante procedure algoritmiche. La discrezionalità giudiziale si configura come il risultato di un'attività valutativa che il giudice esercita nella concretezza del caso sottoposto al suo esame, alla luce delle peculiarità della fattispecie e del contesto nel quale essa si inserisce. In tale processo decisionale assumono inevitabile rilievo la formazione, l'esperienza professionale e la sensibilità interpretativa del singolo magistrato, elementi che contribuiscono a rendere l'esercizio della funzione giurisdizionale difficilmente riconducibile a schemi decisionali rigidamente automatizzati. Occorre, inoltre, guardarsi da indebite generalizzazioni. Non tutti i processi presentano caratteristiche omogenee: processo e procedimento costituiscono categorie concettuali distinte e non sovrapponibili e, soprattutto, le diverse branche dell'ordinamento rispondono a finalità differenti e si fondano su logiche proprie. In questa prospettiva, il processo amministrativo presenta peculiarità strutturali e funzionali che lo distinguono profondamente tanto dal processo civile quanto da quello penale. Proprio tali specificità rendono particolarmente problematica qualsiasi pretesa di uniformare, attraverso strumenti algoritmici, dinamiche decisionali eterogenee, che richiedono invece valutazioni contestuali e apprezzamenti difficilmente riducibili a modelli computazionali. Analoga considerazione può essere svolta con riferimento all'\emph{overruling}. La possibilità di ricondurre tale fenomeno a procedure integralmente realizzabili mediante strumenti di intelligenza artificiale appare, infatti, altamente problematica, oltre che di dubbia utilità pratica, trattandosi di uno degli istituti più strettamente connessi alla dimensione evolutiva e propriamente umana dell'esperienza giuridica. L'\emph{overruling} costituisce, per sua natura, il risultato di un processo interpretativo che si sviluppa nel tempo e che trova alimento nelle trasformazioni della società, dei valori collettivi e della coscienza giuridica. Sotto questo profilo, i giudici si trovano spesso in una posizione privilegiata per cogliere i mutamenti sociali emergenti, poiché operano a diretto contatto con la realtà concreta e con le istanze che da essa provengono. Non di rado, pertanto, l'evoluzione giurisprudenziale, soprattutto in ambito amministrativistico, precede l'intervento del legislatore, il quale è chiamato a osservare e disciplinare fenomeni sociali già in corso di trasformazione\footnote{La tradizione del diritto amministrativo generale italiano mostra come, non di rado, sia stato il legislatore a seguire l'elaborazione giurisprudenziale piuttosto che il contrario, legittimando ex post quella che è stata efficacemente definita ``normazione giurisprudenziale'': cfr., tra gli altri, M. Mazzamuto, \emph{The Formation of the Italian Administrative Justice System, European Common Principles of Administrative Law and ``Jurisdictionalization'' of Administrative Justice in the Nineteenth Century}, in G. della Cananea, S. Mannoni (a cura di), \emph{Administrative Justice Fin de siècle}, Oxford, 2021, p.~282 ss.; M. Clarich, \emph{I materiali della legge nel giudizio amministrativo}, in \emph{Dir. amm.}, 2022, p.~953 ss.}. L'attività interpretativa che conduce al superamento di orientamenti consolidati presuppone una capacità di comprendere il contesto storico, culturale ed etico di riferimento, individuando i segnali di cambiamento che attraversano la società. In tale prospettiva, l'essere umano, in quanto soggetto inserito nelle dinamiche relazionali e sociali della comunità cui appartiene, conserva una funzione difficilmente sostituibile da modelli computazionali fondati sull'elaborazione di dati e precedenti. In questa prospettiva, l'introduzione di sistemi di IA giudicante, soprattutto nell'ambito delle decisioni caratterizzate da elevati margini di discrezionalità o suscettibili di determinare fenomeni di \emph{overruling}, rischierebbe di rappresentare l'ennesima manifestazione delle tendenze semplificatrici proprie dello Stato modernista\footnote{Si segnala per una precisa analisi delle contraddizioni del modernismo F. Bertoni, \emph{La faccia oscura del modernismo}, in \emph{Frontiere/fratture. Il senso del tempo nella storia della letteratura italiana}, Palermo, 2024, pp.~83-103.}. La pretesa di rendere pienamente leggibile, prevedibile e formalizzabile l'attività ermeneutica svolta dal giudice sottende, infatti, una concezione della funzione giurisdizionale che tende a ridurne la complessità a una sequenza di operazioni suscettibili di essere tradotte in modelli computazionali. Una simile impostazione sembra orientata, più o meno consapevolmente, verso una progressiva standardizzazione della produzione giurisdizionale, nella quale il provvedimento giudiziario verrebbe assimilato a un prodotto seriale, generato secondo procedure uniformi e replicabili. Tuttavia, l'attività del giudice amministrativo spesso travalica la mera interpretazione della norma, implicando valutazioni complesse e un'attenta ponderazione tra interessi pubblici e privati eterogenei, nonché l'esercizio di giudizi di valore e la considerazione delle peculiarità del caso concreto. Tali attività presuppongono una comprensione del contesto sociale, culturale e istituzionale nel quale la controversia si colloca, elemento che consente di cogliere la dimensione intrinsecamente contestuale dell'attività giurisdizionale\footnote{Come osservano A. Garapon, J. Lessègue, \emph{La giustizia digitale}, cit., p.~275. A giustificare la riserva di umanità nell'esercizio della funzione giurisdizionale è quella ``saggezza pratica'', la \emph{phronesis} aristotelica, intesa quale capacità di promettere e di impegnarsi a trasformare sé stessi.}. La caratteristica essenziale del giudicante risiede, forse, proprio nel suo essere parte della comunità che è chiamato a regolare mediante l'esercizio della funzione giurisdizionale. In questo senso, il giudicante continua a incarnare l'idea aristotelica dell'uomo quale \emph{ζῷον πολιτικόν}\footnote{G. Angelini, \emph{L'uomo come ζῷον πολιτικόν. Un'ipotesi interpretativa di un lemma fondamentale del pensiero aristotelico}, in \emph{Scienza \& Politica. Per Una Storia Delle Dottrine}, 30, 2018, p.~131 ss.}, soggetto immerso nelle relazioni sociali, nei valori condivisi e nelle trasformazioni della comunità di riferimento. Alla luce delle considerazioni sin qui svolte, appare opportuno abbandonare qualsiasi prospettiva che attribuisca all'intelligenza artificiale una funzione sostitutiva del giudice amministrativo, dovendosi piuttosto valorizzarne il ruolo di strumento ausiliario a supporto dell'attività giurisdizionale. La questione realmente rilevante non consiste nello stabilire se un sistema algoritmico possa assumere integralmente le funzioni del magistrato, bensì nel comprendere in quale misura esso possa contribuire a supportarne l'attività decisionale, migliorando l'efficienza e la qualità del processo senza comprometterne le componenti propriamente umane. In tale prospettiva, l'algoritmo non si colloca come un \emph{tertium genus} tra fatto e diritto, né assume una funzione autonoma nel processo di formazione della decisione. Esso si limita, piuttosto, a operare quale strumento di elaborazione, organizzazione e analisi dei dati, funzionale ad accrescere l'efficienza dell'attività giudiziaria senza incidere sul nucleo propriamente valutativo e decisorio della funzione giurisdizionale, che rimane integralmente riservato al giudice\footnote{B. Custers, E. Fosch-Villaronga, \emph{Law and Artificial Intelligence: Regulating AI and Applying AI in Legal Practice}, cit., spec. cap. I--III.}. Una simile impostazione appare preferibile poiché consente di valorizzare le potenzialità dell'intelligenza artificiale senza smarrire la centralità dell'interprete umano. Essa si mantiene, inoltre, maggiormente distante da uno dei rischi più insidiosi che caratterizzano le moderne forme di governo della complessità sociale: la tendenza a osservare la realtà secondo la logica del \emph{seeing like a State}\footnote{J. Scott, \emph{Seeing like the State}, cit.}, vale a dire attraverso una prospettiva verticale e semplificatrice che, per esigenze di leggibilità e controllabilità, riduce la ricchezza delle interazioni umane e istituzionali a schemi astratti, uniformi e facilmente classificabili. Il pericolo insito in tale approccio risiede nella pretesa di ricondurre fenomeni intrinsecamente eterogenei entro categorie uniformi e standardizzate, sacrificando quelle peculiarità fattuali, sociali e culturali che spesso costituiscono il nucleo più rilevante dell'attività interpretativa e dell'esercizio della discrezionalità giudiziale. L'impiego dell'intelligenza artificiale quale mero strumento ausiliario del giudice consente, invece, di coniugare i vantaggi derivanti dall'automazione e dall'elevata capacità di elaborazione dei dati con la salvaguardia della sensibilità interpretativa e valutativa che costituisce una componente essenziale della decisione giurisdizionale.

\end{document}
